%------------------------------------------------------------------------------%
% FUNCTIONAL ANALYSIS AND APPLICATIONS                                         %
%------------------------------------------------------------------------------%
% File  : Functional_Analysis_and_Applications.tex                             %
% Author: Klaus Hoeflinger, Beethovenstrasse 11, D-73117 Wangen                %
% Email : khoeflinger@t-online.de                                              %
%------------------------------------------------------------------------------%

%------------------------------------------------------------------------------%
%                       DOCUMENT CLASS and INCLUDE FILES                       %
%------------------------------------------------------------------------------%
\documentclass[11pt,twoside]{article}
\usepackage{geometry}
\geometry{paperheight=241mm,
          paperwidth=152mm,
          top=22mm,
          bottom=17mm,
          left=20mm,
          right=20mm,
          textwidth=112mm,
          textheight=202mm,
          headheight=10mm,
          headsep=3mm,
          footskip=9mm,
          includehead,
          includefoot,
          marginparsep=0mm,
          marginparwidth=0mm}

\renewcommand{\baselinestretch}{0.945}\normalsize

\AtBeginDocument{\setlength\abovedisplayskip{2mm}}
\AtBeginDocument{\setlength\belowdisplayskip{2mm}}

%------------------------------------------------------------------------------%
% define title formats                                                         %
%------------------------------------------------------------------------------%
\usepackage{titlesec}

\renewcommand{\thesection}{\arabic{section}}
\renewcommand{\sfdefault}{FiraSans}
\renewcommand{\ttdefault}{pcr}
\renewcommand{\rmdefault}{antiqua}

\titleformat{\part}[display]
{\normalfont \large \rmfamily \bfseries \centering}
{\small{\thepart.}}{2pt}{}

\titleformat{\section}[runin]
{\normalfont \rmfamily \bfseries}
{\thesection}{1pt}{.\;}[.]

%\titleformat{\section}
%{\normalfont \bfseries \small \filcenter}
%{\thesection.\;}{0mm}{}

\titleformat{\subsection}[runin]
{\normalfont \itshape}
{}{0mm}{}

\titleformat{\subsubsection}[runin]
{\normalfont \itshape}
{}{}{}{}

\titleformat{\paragraph}[runin]
{\normalfont}
{}{}{}{}

\titleformat{\subparagraph}[runin]
{\normalfont}
{}{}{}{}

\titlespacing*{\part}          { 0pt}{ 0pt}{ 8pt}
\titlespacing*{\section}       { 0pt}{ 7pt}{ 4pt}
\titlespacing*{\subsection}    { 0pt}{14pt}{ 6pt}
\titlespacing*{\subsubsection} { 0pt}{ 6pt}{12pt}
\titlespacing*{\paragraph}     { 0pt}{ 6pt}{12pt}
\titlespacing*{\subparagraph}  { 0pt}{ 0pt}{12pt}

%------------------------------------------------------------------------------%
% define enumerate parameters                                                  %
%------------------------------------------------------------------------------%
\usepackage{enumitem}
\setlength{\parindent}{3pt}
\setlength{\parskip}  {0pt}

%\setlist{noitemsep}

\setlist[enumerate]{
         leftmargin=12pt,
         rightmargin=0pt,
         itemsep=1pt,
         itemindent=3pt,
         topsep=3pt,
         labelsep=4pt,
         partopsep=0pt,
         labelwidth=0pt,
         listparindent=0pt,
         parsep=0pt}

\setlist[itemize]{
         leftmargin=12pt,
         rightmargin=0pt,
         itemsep=0pt,
         itemindent=1pt,
         topsep=1pt,
         labelsep=4pt,
         partopsep=1pt,
         labelwidth=0pt,
         listparindent=0pt,
         parsep=0pt}

%------------------------------------------------------------------------------%
% define packages                                                              %
%------------------------------------------------------------------------------%
\usepackage{amssymb}
\usepackage{slantsc}
\usepackage{amsmath}
\usepackage{amsthm}
\usepackage{thmtools}
\usepackage{amscd}
\usepackage{verbatim}
\usepackage{latexsym}
\usepackage{multicol}
\usepackage{graphicx}
\usepackage{setspace}
\usepackage[ngerman,english]{babel}
\usepackage[protrusion=true,expansion=true]{microtype}
\usepackage{tikz}
\usetikzlibrary{arrows,arrows.meta,decorations.markings}
\usepackage{mathrsfs}
\usepackage{amsfonts}
\usepackage{datetime2}
\usepackage{textgreek}
\usepackage{upref}
\usepackage{hyperref}
\usepackage{pifont}
\usepackage{relsize}
\usepackage[UKenglish]{isodate}
\usepackage[skip=0mm]{parskip}
\usetikzlibrary{decorations.pathmorphing}
\relsize{-0.05}

\usepackage{mlmodern}
\usepackage[T5,T1]{fontenc}
\usepackage{ragged2e}

%\usepackage{mathabx}
%\usepackage{times}

%\usepackage{fontspec}
%\setmainfont{Nimbus Roman No9 L}
%\usepackage[T1]{fontenc}

%------------------------------------------------------------------------------%
% define page layout                                                           %
%------------------------------------------------------------------------------%
\usepackage{fancyhdr}
\pagestyle{fancy}
\setlength{\headheight}{30.0pt}
\renewcommand{\headrulewidth}{0pt}

\AtBeginDocument{
  \addtolength\abovedisplayskip{-0.0\baselineskip}
  \addtolength\belowdisplayskip{-0.0\baselineskip}
}

\fancyhead{}
\fancyfoot{}
\addtolength{\headwidth}{0.02\marginparwidth}

\makeatletter
\let\ps@plain\ps@fancy
\makeatother

\renewcommand\sectionmark[1]
{
 \def\sectionname{#1}
 \markright{\thesection #1}
}

\makeatletter
\@addtoreset{section}{part}
\makeatother

%------------------------------------------------------------------------------%
% define footnote without number                                               %
%------------------------------------------------------------------------------%
\newcommand{\myfootnote}[1]{%
\renewcommand{\thefootnote}{}%
\footnotetext{#1}%
\renewcommand{\thefootnote}{\arabic{footnote}}%
}

%------------------------------------------------------------------------------%
% define numbering in equations                                                %
%------------------------------------------------------------------------------%
\numberwithin{equation}{section}

%------------------------------------------------------------------------------%
% define newtheorem                                                            %
%------------------------------------------------------------------------------%
\declaretheoremstyle[
headfont=\scshape, 
headindent = 0mm, %\parindent,
%postheadhook = {\hspace{0mm}\newline},
%postheadspace = \newline, 
spaceabove = 3mm, 
spacebelow = 3mm,
numberwithin=section,
name=Theorem]{khMATH}
\declaretheorem[style=khMATH]{theorem}

\declaretheoremstyle[
headfont=\bfseries, 
headindent = 0mm, %\parindent,
%postheadhook = {\hspace{0mm}\newline},
%postheadspace = \newline, 
spaceabove = 3mm, 
spacebelow = 3mm,
numberwithin=part,
name=Theorem]{khMATH1}
\declaretheorem[style=khMATH1]{theorem1}

\declaretheoremstyle[
headfont=\scshape, 
headindent = 0mm, %\parindent,
%postheadhook = {\hspace{0mm}\newline},
%postheadspace = \newline, 
spaceabove = 3mm, 
spacebelow = 3mm,
numberwithin=section,
name=Corollary]{khMATH1}
\declaretheorem[style=khMATH1]{corollary}

%            name         counter     label              numeration-mode
\theoremstyle{plain}
\newtheorem*{introduction}            {\sc Introduction}
\newtheorem {standard}                {}                 [section]
\newtheorem*{definition}              {\sc Definition}
\newtheorem{satz1}                    {\sc Satz}
\newtheorem{lemma}                    {\sc Lemma}
\newtheorem*{proposition}             {\sc Proposition}
%\newtheorem{corollary}               {\sc Corollary}
\newtheorem*{corollar}                {\sc Korollar}
\newtheorem*{remark}                  {\sc Remark}
\newtheorem*{remarks}                 {\sc Remarks}
\newtheorem*{example}                 {\sc Example}
\newtheorem*{examples}                {\sc Examples}
\newtheorem*{theo}                    {\sc Theorem}

%------------------------------------------------------------------------------%
% define tabular                                                               %
%------------------------------------------------------------------------------%
\usepackage{tabularx}
\newcolumntype{L}[1]{>{\raggedright\arraybackslash}p{#1}}
\newcolumntype{C}[1]{>{\centering\arraybackslash}  p{#1}} 
\newcolumntype{R}[1]{>{\raggedleft\arraybackslash} p{#1}} 

%------------------------------------------------------------------------------%
% define \sh, \zB                                                              %
%------------------------------------------------------------------------------%
\def\dsh{{d.\,h.}}
\def\ie{{i.\,e.}}
\def\zB{z.\,B.}
\renewcommand{\abstractname}{ }

%------------------------------------------------------------------------------%
% change default spacing for cases                                             %
%------------------------------------------------------------------------------%
\makeatletter
\renewcommand*\env@cases[1][1.6]{%
  \let\@ifnextchar\new@ifnextchar
  \left\lbrace
  \def\arraystretch{#1}%
  \array{@{}l@{\quad}l@{}}%
}
\makeatother

\newenvironment{rcases}
  {\left.\begin{aligned}}
  {\end{aligned}\right\rbrace}

%------------------------------------------------------------------------------%
% change default for \predisplaypenalty                                        %
%------------------------------------------------------------------------------%
\predisplaypenalty=990
\allowdisplaybreaks \numberwithin{equation}{section}

%------------------------------------------------------------------------------%
% general notations                                                            %
%------------------------------------------------------------------------------%
\def\*{\star}
\def\={\,=\,}
\def\+{\,+\,}
\def\xsubset{{\,\subset\,}}
\def\xtimes{{\,\times\,}}
\def\xeof{{\,\in\,}}
\def\eq{{\,=\,}}
\def\xto{{\,\to\,}}
\def\eqdef{\,:=\,}
\def\:{{\,:\,}}
\def\ele{{\,\in\,}}
\def\exists{\mbox{ exists }}
\def\and{\mbox{ and }}
\def\for{\mbox{ for }}
\def\fa{\mbox{ for all }}
\def\fe{\mbox{ for each }}
\def\qfa{\quad \mbox{ for all }}
\def\df{\,\dotfill}
\def\iff{if and only if }
\def\wrt{with respect to}

\makeatletter
\newcommand \Df {\leavevmode \cleaders \hb@xt@ .70em{\hss .\hss }\hfill \kern \z@}
\makeatother

\def\sql{\mathfrak{a}}               % sesquilinear form a
\def\DF{B}                           % Dirichlet form a
%------------------------------------------------------------------------------%
% number sets and special elements                                             %
%------------------------------------------------------------------------------%
\def\P{\mathbb{H}}                   % half plane of $\C$
\def\J{I}                            % interval I
\def\N{{\mathbb{N}}}                 % unsigned integers
\def\Q{{\mathbb{Q}}}                 % rational integers
\def\Z{{\mathbb{Z}}}                 % integers
\def\R{{\mathbb{R}}}                 % real numbers
\def\Rn{{\mathbb{R}}^n}              % real numbers in Rn
\def\Rd{{\mathbb{R}}^d}              % real numbers in Rn
\def\C{{\mathbb{C}}}                 % complex numbers
\def\F{{\mathbb{K}}}                 % arbitrary field
\def\H{{\mathbb{H}}}                 % half of the complex plane
\def\Torus{{\mathbb{T}}}             % complex circle line
\def\T{\tau}                         % upper bound of interval (0,t)
\def\sign#1{\operatorname{sign}(#1)} % sign of a number

%------------------------------------------------------------------------------%
% set theory                                                                   %
%------------------------------------------------------------------------------%
\def\G{G}                            % domain $G$
\def\clG{\overline{\G}}              % closure of $G$
\def\bndG{\partial \G}               % boundary of $G$
\def\setA{\mathcal{A}}               % set A
\def\setB{\mathcal{B}}               % set B
\def\setC{\mathcal{C}}               % set C
\def\setM{M}                         % set M
\def\setN{N}                         % set N
\def\setS{\mathcal{S}}               % set S
\def\famF{\mathfrak{F}}              % family of sets
\def\indI{\mathcal{I}}               % index set I
\def\pot#1{\mathfrak{P}(#1)}         % power set

\def\relR{\mathbf{R}}                % relation R
\def\mapf{f}                         % mapping f
\def\mapg{g}                         % mapping g

\def\set#1#2{\{ \, #1 \,: \, #2 \, \}}
\def\tensor#1#2{#1 \otimes #2}
\def\Tens{\oplus}

\def\cal#1{{\mathcal{#1}}}
\def\aut#1{\operatorname{Aut}(#1)}

\def\cross#1#2{#1 {\,\times\,} #2}
\def\multicross#1#2{#1 \times  \ldots \times  #2}
\def\multitens#1#2 {#1 \otimes \ldots \otimes #2}

%------------------------------------------------------------------------------%
% group theory                                                                 %
%------------------------------------------------------------------------------%
\def\1{{\mathsf{1}}}                % unit element of a monoid
\def\0{{\mathsf{0}}}                % unit element of an Abelian monoid
\def\kernel#1{\mathfrak{N}(#1)}     % kernel of a homomorphism
\def\cokernel#1{\mathfrak{C}{(#1)}} % cokernel of a homomorphism
\def\Hom#1{\operatorname{Hom}(#1)}

%------------------------------------------------------------------------------%
% linear algebra                                                               %
%------------------------------------------------------------------------------%
\def\vecV{\mathit{V}}               % vector space V
\def\vecH{\mathit{H}}               % vector space H
\def\vecW{\mathit{W}}               % vector space W
\def\vecU{\mathit{U}}               % subspace U
\def\convC{\mathcal{C}}             % convex set C
\def\basH{\mathfrak{B}}             % Hamel base B

\def\LO#1{\mathscr{L}(#1) }         % algebra of linear transformations

\def\scalar#1#2{ \langle #1,#2 \rangle }
\def\trace#1{\operatorname{tr}(#1)}
\def\dim#1{{\operatorname{dim} \, #1 }}
\def\codim#1{{\operatorname{codim} \, #1 }}
\def\dim{\operatorname{dim}}
\def\codim{\operatorname{codim}}
\def\ind{\operatorname{ind}}

\def\genG{\cal{G}}

\def\algA{\cal{A}}
\def\algB{\cal{B}}
\def\algU{\cal{U}}

\def\idI{\cal{I}}

%------------------------------------------------------------------------------%
% topology                                                                     %
%------------------------------------------------------------------------------%
\def\compK{{K}}          % compact set C
\def\openO{\mathcal{O}}  % open set O
\def\openU{\mathcal{U}}  % open set U
\def\U{\mathcal{U}}      % neighbourhood U
\def\V{\mathcal{V}}      % neighbourhood V
\def\d{\mathit{d}}       % metric function d

\def\interior#1{\mathring{#1}}
\def\closure#1{{\overline{#1}}}

\def\topT{\mathcal{T}}   % topology T
\def\topX{X}             % topological space X
\def\topY{Y}             % topological space Y
\def\topZ{Z}             % topological space Z
\def\metrS{M}            % metric space M

\def\grpG{G}             % topological group G
\def\grpA{A}             % topological group A
\def\grpB{B}             % topological group B
\def\grpC{C}             % topological group C
\def\grpH{H}             % topological group H
\def\grpU{U}             % topological group U
\def\grpV{V}             % topological group V
\def\topG{G}             % topological group G
\def\topH{H}             % topological group H
\def\topU{U}             % topological subgroup U

\def\subS{\cal{S}}
\def\riR{R}              % ring R

%------------------------------------------------------------------------------%
% calculus                                                                     %
%------------------------------------------------------------------------------%
\def\supp#1{\operatorname{supp}(#1)}
\def\singsupp#1{\operatorname{sing \, supp}(#1)}

\def\re{\operatorname{Re}}
\def\im{\operatorname{Im}}
\def\grad{\operatorname{grad}}

%\def\abs#1 {\left|#1\right|}
\def\abs#1 {\lvert #1 \rvert}
%\def\abs#1 {\left|\,#1\,\right|}

%------------------------------------------------------------------------------%
% measure theory                                                               %
%------------------------------------------------------------------------------%
\def\salgA{\Sigma}               % sigma-algebra S
\def\salgB{\cal{B}}              % sigma-algebra B
\def\Borel#1{\cal{B}_{#1}}       % Borel-algebra 
\def\LB{\lambda}                 % Lebesgue-Borel measure

%------------------------------------------------------------------------------%
% function spaces                                                              %
%------------------------------------------------------------------------------%
\def\Lp{L^p(\G)}                 % spaces L_p
\def\Lq{L^q(\G)}                 % spaces L_q
\def\Lh{L^2(\G)}                 % spaces L_2
\def\Li{L^{\infty}(\G)}          % spaces L_infty
\def\Lc{L_{loc}^1(\G)}           % spaces L_1^{loc}
\def\Ck{C^k(\G)}                 % spaces C_k

\def\BV#1{BV(#1) }               % set of functions of bounded variation
\def\AC#1{AC(#1) }               % set of absolutely continuous functions

\def\MP#1{\mathscr{M}^+(#1) }    % set of positive measures
\def\MS#1{\mathscr{M}^-(#1) }    % vector space of finite signed measures
\def\MC#1{\mathscr{M}(#1) }      % vector space of complex measures
\def\MCR#1{\mathscr{M}_{r}(#1) } % vector space of regular complex measures

\def\SobW{W^{k,p}(\G)}                % Sobolev space W
\def\SobO{\mathring{W}^{k,p}(\G)}     % Sobolev space W
%\def\WN#1#2{\mathring{W}^{#1}(#2) }   % Sobolev space W
\def\WN#1#2{W_0^{#1}(#2) }   % Sobolev space W
%\def\HN#1#2{\mathring{H}^{#1}(#2) }   % Sobolev space H
\def\HN#1#2{H_0^{#1}(#2) }   % Sobolev space H
%------------------------------------------------------------------------------%
% functional analysis                                                          %
%------------------------------------------------------------------------------%
\def\snorm#1{\lVert #1 \rVert}
\newcommand{\norm}[1]{\left\lVert #1 \right\rVert}

\def\topV{V}                       % topological vector space V
\def\topW{W}                       % topological vector space W
\def\normS{X}                      % normed space

\def\banX{\mathit{X}}              % Banach space X
\def\banY{\mathit{Y}}              % Banach space Y
\def\banZ{\mathit{Z}}              % Banach space Z
\def\dbanX{\mathit{X^*}}           % dual space of Banach space X
\def\dbanY{\mathit{Y^*}}           % dual space of Banach space Y
\def\dbanZ{\mathit{Z^*}}           % dual space of Banach space Z

\def\banA{\mathit{A}}              % Banach algebra A
\def\banB{\mathit{B}}              % Banach algebra B
\def\banC{\mathit{C}}              % C*-algebra C


\def\domain#1{\mathfrak{D}(#1)}    % domain
\def\range#1{\mathfrak{R}(#1)}     % range

\def\isoJ{\mathfrak{J}}            % isometry J
\def\repA{{\cal{A}_{\sql}}}        % representation operator A
\def\linS{{S}}                     % linear operator S
\def\linG{{G}}                     % linear operator G
\def\linL{{L}}                     % linear operator L
\def\linT{{T}}                     % linear operator A
\def\linA{\mathit{A}}              % linear operator A
\def\linB{{B}}                     % linear operator B
\def\linM{{M}}                     % linear operator M
\def\linU{{U}}                     % linear operator U
\def\linI{{I}}                     % linear operator I
\def\A{\cal{A}}                    % linear differential operator A
\def\BDO{\cal{B}}                  % linear boundary differential operator B
\def\linU{{U}}                     % unitary operator U
\def\linP{{P}}                     % projection P
\def\compA{k}                      % compact operator k
\def\linFR{f}                      % finite rank operator f
\def\fredA{F}                      % Fredholm operator F

\def\HAM{\cal{H}}                  % Hamilton operator
\def\PA{\partial^{\alpha}}         % elementary differential operator
\def\DO{\cal{A}}                   % linear differential operator
\def\BC{\cal{B}}                   % boundary operator
\def\SLO{\cal{L}_{p,q}}            % Sturm-Liouville operator
\def\MO{\cal{L}_{M}}               % momentum operator

\def\HS#1{S(#1)           }        % algebra of Hilbert-Schmidt operators
\def\FR#1{F(#1)           }        % algebra of finite rank operators
\def\TC#1{N(#1)           }        % algebra of trace class operators
\def\BU#1{\mathscr{U}(#1) }        % algebra of unitary linear operators
\def\BO#1{\mathscr{L}(#1) }        % algebra of bounded linear operators
\def\CO#1{\mathscr{C}(#1) }        % algebra of compact linear operators
\def\CL#1{\mathscr{C}(#1) }        % algebra of compact linear operators
\def\FO#1{\Phi(#1) }               % set of Fredholm operators
\def\FK#1#2{\Phi_{#2}(#1) }        % set of Fredholm operators with index k

\def\hilH{\mathit{H}}              % Hilbert space H
\def\clU{\mathit{U}}               % closed subspace U

\def\orthO{\mathfrak{O}}           % orthonormal system O
\def\orthS{\mathfrak{S}}           % orthonormal system S
\def\basB{\mathfrak{B}}            % basis B of a Hilbert space

\def\GL#1 {\mathbf{GL}(#1)}        % linear group
\def\OG#1 {\mathbf{O}(#1)}         % linear group
\def\SOG#1 {\mathbf{SO}(#1)}       % linear group

%------------------------------------------------------------------------------%
% semigroup theory                                                             %
%------------------------------------------------------------------------------%
\def\SG{$\{S(t)\}_{t \geq 0}$}
\def\TG{$\{T(t)\}_{t \geq 0}$}
\def\GG{$\{G(t)\}_{t \geq 0}$}
\def\EG{$\{\cal{E}(t)\}_{t \geq 0}$}

%------------------------------------------------------------------------------%
% distribution theory                                                          %
%------------------------------------------------------------------------------%
\def\disF{F}                  % distribution F
\def\disG{G}                  % distribution G
\def\disH{H}                  % distribution H
\def\disU{U}                  % distribution U
\def\disT{T}                  % distribution T
\def\SF{C^{\infty}(\G)}       % space of smooth functions
\def\TF{C_c^{\infty}(\G)}     % space of compactly supported smooth functions
\def\TFRd{C_c^{\infty}(\R^d)} % space of test functions
\def\SRd{\mathscr{S}(\R^d)}   % Schwartz space
\def\SR1{\mathscr{S}(\R)}     % Schwartz space
\def\B1#1{C^1(#1) }           % space of continuously differentiable functions
\def\Bm#1{C^m(#1) }           % space of continuously differentiable functions
\def\Ck{C^k(\G)}
\def\TD{\mathscr{S}'(\R^d)}   % space of temperated distributions
\def\E{\cal{E}}               % space of distributions with compact support
\def\D{\cal{D}}               % D
\def\FT{\mathcal{F}}          % Fourier transformation F
\def\FT{\mathscr{F}}          % Fourier transformation F

%------------------------------------------------------------------------------%
% other definitions                                                            %
%------------------------------------------------------------------------------%
\def\Proof{\textsc{Proof}}
\def\FRECHET{Fr\'{e}chet}
\def\ARZELA{Arzel\'{a}}
\def\GARDING{G\r{a}rding}
\def\HOELDER{H\"older}
\def\SCHROEDINGER{Schr\"odinger}
\def\JOERGENS{J\"orgens}
\def\WUEST{W\"uest}
\def\KAEHLER{K\"aehler}
\def\HOEFLINGER{H\"oflinger}
\def\POINCARE{Poincar\'{e}}
\def\FA{Functional Analysis}

\def\Author   {K.\,{\HOEFLINGER}}
\def\AuthorUC {K.\,HOEFLINGER}
\def\Version  {1.0}
\def\Email    {khoeflinger@t-online.de}
\def\khMath   {https://khoefli.github.io/khMath}
\def\Year     {2026}


\usepackage{newunicodechar}
\newunicodechar{²}{\ifmmode{}^2\else\texttwosuperior\fi}

\renewcommand\thesection{\arabic{section}}

\titleformat{\part}[display]
{\normalfont \large \rmfamily \bfseries \centering}
{\small{\thepart.}}{2pt}{}

\titleformat{\section}%[runin]
{\normalfont \bfseries \filcenter}
{\S \thesection.\;}{0mm}{}

\titlespacing*{\section}   {0pt}{16pt}{ 4pt}
\titlespacing*{\subsection}{0pt}{10pt}{12pt}

\newlength{\bibitemsep}\setlength{\bibitemsep}{.5\baselineskip plus .05\baselineskip minus .05\baselineskip}
\newlength{\bibparskip}\setlength{\bibparskip}{0pt}
\let\oldthebibliography\thebibliography
\renewcommand\thebibliography[1]{%
  \oldthebibliography{#1}%
  \setlength{\parskip}{\bibitemsep}%
  \setlength{\itemsep}{\bibparskip}%
}

%------------------------------------------------------------------------------%
%                                BEGIN OF DOCUMENT                             %
%------------------------------------------------------------------------------%
\begin{document}
\thispagestyle{empty}

\def\Title{Functional Analysis and Applications}

\centerline{\large{\textbf{\Title}}}
\vspace{4pt}
\centerline{\small{\textsc{\Author}}}
\vspace{10pt}

\myfootnote { Latest Version: \today.}
\myfootnote { Email: \textit{{\Email}}, Url: \href{\khMath}{\khMath}.}


\fancyhead[OL]{\small{\thepage}}
\fancyhead[ER]{\small{\thepage}}

%------------------------------------------------------------------------------%
%                                 INTRODUCTION                                 %
%------------------------------------------------------------------------------%
This text consists of a series of articles
on the mathematical area of {functional analysis}
and of a final one concerning with the applications of the abstract theory
to linear partial differential equations,
which arise as models for problems
in mathematical physics and engineering science.

\subparagraph{}
Functional analysis was mainly developed
in the first half of the 20$^{\textrm{th}}$ century,
and valuable contributions have been done by
{D.\;Hilbert},
{S.\;Banach},
{F.\;Riesz},
{J.\;von Neumann},
{F.\,E.\;Browder},
{L.\;Schwartz},
{J.-L.\;Lions} and others.

\subparagraph{}
The topics within functional analysis are mainly divided into two parts:
the linear and the non-linear theory.
In our treatise we shall exclusively deal with the linear one,
hence the fundamental concept of {"linear operator"}
will play a central role in our considerations.
Especially
what historically is meant by the concept of unbounded operator
is of great importance,
since the theory of this specific class of mappings can be successfully applied
to the study of partial differential operators
and their related boundary and initial-value problems.

\paragraph{}
It is assumed that
the reader is familiar with elementary concepts in linear algebra,
general topology and measure theory.
In fact,
functional analysis can be seen as the synthesis of the both first subjects,
especially {\wrt} infinite-dimensional linear spaces.
Further preliminaries
will be introduced as soon as there is the need to use them.

\subparagraph{}
It is not the idea of this text to give all the proofs of the stated theorems,
except in the case where these are elementary.
Instead,
the main purpose is a transparent and overall presentation of the topics,
especially for readers, which have learned functional analysis in the past
and want to have a refresh of the main results now.

\paragraph{} %- honoration to Prof. Dr. Werner
I wrote the treatise in honour of Prof.\;Dr.\;P.\;Werner (1932-2018),
who conveyed my fascination of these beautiful mathematical topics
by his lessons hold in 1984 and 1985 at
Mathematical Institute A of University Stuttgart.

%------------------------------------------------------------------------------%
%                              TABLE OF CONTENTS                               %
%------------------------------------------------------------------------------%
\newpage
\thispagestyle{empty}
\section*{Table of Contents}
\vspace{20pt}
%------------------------------------------------------------------------------%

\begin{small}
\vspace{7pt}
\centerline{{\footnotesize{Part I}}}
\vspace{1pt}
\centerline{\textbf{Topological and Normed Linear Spaces}}
\vspace{2pt}

\centerline
{
\begin{tabular}{R{5mm}L{93mm}R{5mm}}
%------------------------------------------------------------------------------%
 1. & Generalities on Topological Linear Spaces                      \Df &  1 \\
 2. & Seminorms and Locally Convex Spaces                            \Df &  3 \\
 3. & Normed Linear Spaces                                           \Df &  4 \\
 4. & Strictly and Uniformly Convex Spaces                           \Df &  5 \\
 5. & Banach and Hilbert Spaces                                      \Df &  6 \\
 6. & Spaces of Scalar-Valued Functions and Measures                 \Df &  9 \\
 7. & Hierarchy of Topological and Normed Linear Spaces              \Df & 18 \\
%------------------------------------------------------------------------------%
\end{tabular}
}

\vspace{7pt}
\centerline{{\footnotesize{Part II}}}
\vspace{1pt}
\centerline{\textbf{Linear Operators in Banach Spaces}}
\vspace{2pt}

\centerline
{
\begin{tabular}{R{5mm}L{93mm}R{5mm}}
%------------------------------------------------------------------------------%
 1. & Fundamentals of Operator Theory                                \Df & 19 \\
 2. & Continuous Linear Operators                                    \Df & 20 \\
 3. & Compact Linear Operators                                       \Df & 22 \\
 4. & Closed and Closable Linear Operators                           \Df & 23 \\
 5. & The Closed Graph Theorem                                       \Df & 25 \\
%------------------------------------------------------------------------------%
\end{tabular}
}

\vspace{7pt}
\centerline{{\footnotesize{Part III}}}
\vspace{1pt}
\centerline{\textbf{Adjoint Spaces and Operators}}
\vspace{2pt}

\centerline
{
\begin{tabular}{R{5mm}L{93mm}R{5mm}}
%------------------------------------------------------------------------------%
 1. & Continuous Linear Functionals                                  \Df & 34 \\
 2. & Weak Topologies, Bidual Spaces and Reflexivity                 \Df & 38 \\
 3. & Representation Theorems for Adjoint Spaces                     \Df & 39 \\
 4. & Spaces of Distributions                                        \Df & 52 \\
 5. & Adjoint Operators in Locally Convex Spaces                     \Df & 39 \\
 6. & Adjoint Operators in Banach Spaces                             \Df & 39 \\
 7. & Hilbert Adjoints and Normal Operators                          \Df & 40 \\
%------------------------------------------------------------------------------%
\end{tabular}
}

\vspace{7pt}
\centerline{{\footnotesize{Part IV}}}
\vspace{1pt}
\centerline{\textbf{Normally Solvable Operators}}
\vspace{2pt}

\centerline
{
\begin{tabular}{R{5mm}L{93mm}R{5mm}}
%------------------------------------------------------------------------------%
 1. & Banach's Closed Range Theorem                                  \Df & 54 \\
 2. & Operators Bounded Below                                        \Df & 56 \\
 3. & The Fredholm Index                                             \Df & 56 \\
 4. & Perturbations of Semi-Fredholm Operators                       \Df & 58 \\
 5. & Basic Theory on Fredholm Operators                             \Df & 59 \\
 6. & Fredholm's Alternative Theorem                                 \Df & 60 \\
 7. & Concepts of Solvability for Linear Equations                   \Df & 63 \\
%------------------------------------------------------------------------------%
\end{tabular}
}

\newpage
\thispagestyle{empty}
\centerline{{\footnotesize{Part V}}}
\vspace{1pt}
\centerline{\textbf{Calculus of Vector-Valued Functions}}
\vspace{2pt}

\centerline
{
\begin{tabular}{R{5mm}L{93mm}R{5mm}}
%------------------------------------------------------------------------------%
 1. & Spaces of Classical Differentiable Functions                   \Df & 90 \\
 2. & Riemann Integration                                            \Df & 90 \\
 3. & Linear Evolution Equations                                     \Df & 90 \\
 4. & Bochner Integration                                            \Df & 90 \\
 5. & Inhomogeneous Cauchy Problems                                  \Df & 90 \\
 6. & Weak Derivatives and Bochner Sobolev-Spaces                    \Df & 90 \\
 7. & Hilbert-Bochner-Sobolev Spaces                                 \Df & 90 \\
 8. & Holomorphic Vector-Valued Functions                            \Df & 90 \\
%------------------------------------------------------------------------------%
\end{tabular}
}

\vspace{7pt}
\centerline{{\footnotesize{Part VI}}}
\vspace{1pt}
\centerline{\textbf{General Spectral Theory}}
\vspace{2pt}

\centerline
{
\begin{tabular}{R{5mm}L{93mm}R{5mm}}
%------------------------------------------------------------------------------%
 1. & Resolvent Sets and Resolvent Functions                         \Df & 72 \\
 2. & Different Kinds of Spectra                                     \Df & 74 \\
 3. & Riesz-Schauder Spectral Theory                                 \Df & 78 \\
 4. & Linear Operators with Pure Point Spectrum                      \Df & 80 \\
 5. & The Riesz-Dunford Functional Calculus                          \Df & 75 \\
 6. & Essential and Discrete Spectrum                                \Df & 75 \\
%------------------------------------------------------------------------------%
\end{tabular}
}

\vspace{7pt}
\centerline{{\footnotesize{Part VII}}}
\vspace{1pt}
\centerline{\textbf{Commutative Banach Algebras}}
\vspace{2pt}

\centerline
{
\begin{tabular}{R{5mm}L{93mm}R{5mm}}
%------------------------------------------------------------------------------%
 1. & Normed Algebras                                                \Df & 72 \\
 2. & Resolvent and Spectrum of Banach Algebra Elements              \Df & 72 \\
 3. & Gelfand Representation                                         \Df & 90 \\
 4. & Involutional Algebras and C*-Algebras                          \Df & 90 \\
 5. & The Gelfand-Naimark Theorem                                    \Df & 90 \\
 6. & Topological and Locally Compact Abelian Groups                 \Df & 90 \\
 7. & Abstract Fourier Transformation                                \Df & 90 \\
 8. & Fourier Transformation in Euclidean Spaces                     \Df & 90 \\
%------------------------------------------------------------------------------%
\end{tabular}
}

\vspace{7pt}
\centerline{{\footnotesize{Part VIII}}}
\vspace{1pt}
\centerline{\textbf{Linear Operators in Hilbert Space}}
\vspace{2pt}

\centerline
{
\begin{tabular}{R{5mm}L{93mm}R{5mm}}
%------------------------------------------------------------------------------%
 1. & Numerical Range and External Field of Regularity               \Df & 90 \\
 2. & Unitary Operators and Isometries                               \Df & 90 \\
 3. & Symmetric Linear Operators                                     \Df & 95 \\
 4. & Self-Adjoint Operators                                         \Df & 97 \\
 5. & Essentially Self-Adjointness                                   \Df & 99 \\
 6. & The von\,Neumann Extension Theory                              \Df & 98 \\
 7. & Perturbations of Self-Adjoint Operators                        \Df &101 \\
%------------------------------------------------------------------------------%
\end{tabular}
}

\newpage
\thispagestyle{empty}
\centerline{{\footnotesize{Part IX}}}
\vspace{2pt}
\centerline{\textbf{Spectral Theory of Self-Adjoint Operators}}
\vspace{2pt}

\centerline
{
\begin{tabular}{R{5mm}L{93mm}R{5mm}}
%------------------------------------------------------------------------------%
 1. & Spectral Properties of Normal Operators                        \Df &103 \\
 2. & Properties of the Resolvent and the Spectrum                   \Df &105 \\
 3. & The Discrete Spectral Theorem                                  \Df &113 \\
 4. & Projections and Projection-Valued Measures                     \Df &113 \\
 5. & Functional Calculus and Spectral Operators                     \Df &115 \\
 6. & The Spectral Theorem for Self-Adjoint Operators                \Df &117 \\
 7. & Positive Self-Adjoint Operators and Cauchy Problems            \Df &125 \\
%------------------------------------------------------------------------------%
\end{tabular}
}

\vspace{7pt}
\centerline{{\footnotesize{Part X}}}
\vspace{1pt}
\centerline{\textbf{Semigroup Theory and Evolution Equations}}
\vspace{2pt}

\centerline
{
\begin{tabular}{R{5mm}L{93mm}R{5mm}}
%------------------------------------------------------------------------------%
 1. & Strongly Continuous Semigroups                                 \Df &130 \\
 2. & Infinitesimal Generators                                       \Df &140 \\
 3. & Accretive Operators and Contractive Semigroups                 \Df &141 \\
 4. & Sectorial Operators and Analytic Semigroups                    \Df &143 \\
 5. & Skew-Adjoint Operators and Unitary Groups                      \Df &145 \\
 6. & Well-Posedness of Second-Order Cauchy Problems                 \Df &148 \\
 7. & The Cosine Family Approach                                     \Df &149 \\
%------------------------------------------------------------------------------%
\end{tabular}
}

\vspace{7pt}
\centerline{{\footnotesize{Part XI}}}
\vspace{1pt}
\centerline{\textbf{Coercive Sesquilinear Forms}}
\vspace{2pt}

\centerline
{
\begin{tabular}{R{5mm}L{93mm}R{5mm}}
%------------------------------------------------------------------------------%
 1. & Generalities on Sesquilinear Forms                             \Df &160 \\
 2. & Coercive Forms and Variational Equations                       \Df &162 \\
 3. & Triples of Hilbert Spaces                                      \Df &164 \\
 4. & Linear Operators Associated with a H-Elliptic Form             \Df &165 \\
 5. & Sectorial Forms and m-Sectorial Operators                      \Df &167 \\
 6. & The Friedrichs Extension Theorem                               \Df &168 \\
 7. & Variational Evolution Equations                                \Df &169 \\
%------------------------------------------------------------------------------%
\end{tabular}
}

\vspace{7pt}
\centerline{{\footnotesize{Part XII}}}
\vspace{1pt}
\centerline{\textbf{Applications to Differential Operators}}
\vspace{2pt}

\centerline
{
\begin{tabular}{R{5mm}L{93mm}R{5mm}}
%------------------------------------------------------------------------------%
 1. & Free Space and Boundary Value Problems                         \Df &172 \\
 2. & Operators of Elliptic, Parabolic and Hyperbolic Type           \Df &173 \\
 3. & Realizations in Banach Spaces of Functions                     \Df &175 \\
 4. & Closeness and Closability of Differential Operators            \Df &176 \\
 5. & The Fredholm Property of Properly Elliptic Operators           \Df &177 \\
 6. & Self-Adjoint Extensions                                        \Df &178 \\
 7. & Eigenvalue Problems in Bounded Domains                         \Df &179 \\
 8. & Variational Boundary Value Problems                            \Df &180 \\
 9. & Linear Differential Operators in Distributional Spaces         \Df &181 \\
10. & Fourier Transform Methods                                      \Df &182 \\
%------------------------------------------------------------------------------%
\end{tabular}
}
\end{small}

%------------------------------------------------------------------------------%
%                                   CONTENT                                    %
%------------------------------------------------------------------------------%
\pagenumbering{arabic}
\fancyhead[OL]{\small{\thepage}}
\fancyhead[ER]{\small{\thepage}}
\setcounter{page}{0}

\newpage
\thispagestyle{empty}
\fancyhead[C]{\small{I. Topological and Normed Linear Spaces}}
%------------------------------------------------------------------------------%
\part{Topological and Normed Linear Spaces}                                    %
%------------------------------------------------------------------------------%
Our aim in this first part
is to review those fundamental classes of topological linear spaces,
in which the analysis of so-called \textit{linear operators} is usually done.
These classes include,
besides some generic classes like
{locally convex spaces} and {completely metrizable spaces},
especially the class of {normed linear spaces} and its subclasses,
where we want to emphasize {Banach} and {Hilbert spaces}.
In fact,
there are a many normed linear spaces appearing in higher analysis,
whose elements are either scalar-valued functions or measures.

\section{Generalities on Topological Linear Spaces}
%------------------------------------------------------------------------------%
Basically,
the scalar field $\F$ of the linear spaces under consideration in this text
will either be the field $\R$ of real numbers
or the field $\C$ of complex numbers.
Notice that these fields are furnished
with the so-called {standard topology},
whose subbase is given by the totality of open intervals in the real line
and by the totality of open disks in the complex plane,
respectively.
In general,
if results on topological or normed linear spaces
prove to be independent of the choice of the scalar field,
we shall use the symbol $\F$ instead of $\R$ or $\C$. 

\subparagraph{} %- topological linear spaces
By a \textit{topological linear space} is meant an ordinary vector space $\vecV$
over a topological field $\F$ (usually that of real or complex numbers),
together with a topology $\topT_{\,\vecV}$ 
such that addition $(x,y) \mapsto x\,{+}\,y$
and outer multiplication $(\lambda,x) \mapsto \lambda x$ are continuous.
Hence a multiplicity of topological concepts,
like
the \textit{closure} of a set,
the \textit{neighbourhood} of a point,
the \textit{convergence} of a sequence,
the \textit{continuity} of a mapping and others
are meaningful in the context of topological linear spaces.

\paragraph{} %- Hausdorff's separation axiom
We continue with the introduction of a basic property
that our topological linear spaces under consideration have to own:
all singletons $\{x\}$ with $x \xeof \vecV$ are closed sets,
which is equivalent to say
that the topology $\topT_{\,\vecV}$
fulfills \textsc{Hausdorff}'s separation axiom.
In other words,
each pair of points $x,y \xeof \banX$ can be separated
by open sets $\openO_x, \openO_y$,
that is,
there are $\openO_x \subset \vecV$ and $\openO_y \subset \vecV$ with
$x \xeof \openO_x$,
$y \xeof \openO_y$ and
$\openO_x \cap \openO_y \eq \emptyset$.
The validness of this separation axiom implies not only
that every singleton $\{x\}$ is closed,
but also
that the limit of a convergent sequence of elements in $\vecV$ is unique.

\paragraph{} %- metrizable linear spaces
Among a lot of properties,
a topological linear space may have,
there are two special ones
that are of importance for our purposes:

\begin{itemize}[leftmargin=12pt,itemsep=5pt]
\item[1.]
A topological linear space $\vecV$ is said to be \textit{metrizable},
if there is a function
$\d\:\vecV \xtimes \vecV \xto [0,\infty)$
fulfilling the properties of a \textit{metric}\footnote{
\,Given a non-empty set $\metrS$,
a mapping $\d\:\cross{\metrS}{\metrS} \xto \R$ is called a \textit{metric}
on $\metrS$
iff the following four conditions hold for all $x,y,z \in \metrS$:

\begin{itemize}[leftmargin=40pt,itemsep=0pt]
\item[(M-1)]
$d(x,y) \geq 0$,
\item[(M-2)]
$d(x,y) \eq 0$ if and only if $x=y$,
\item[(M-3)]
$d(x,y) \eq d(y,x)$;
\item[(M-4)]
$d(x,z) \leq d(x,y) + d(y,z)$ \hspace{3mm}(\textit{triangle inequality}).
\end{itemize}

Such a function $\d$ gives rise to a topology in $\metrS$,
called the \textit{metric topology} associated with $\d$,
and a subbase is given by all open balls
$B_r(x) \eqdef \{y \xeof \metrS: \d(x,y) < r\}$
with center $x \xeof \metrS$ and radius $r > 0$.
}
such that the metric topology associated with $\d$
resembles the given topology $\topT_{\,\vecV}$.

\item[2.]
A sequence $\{x_n\}_{n \in \N}$ with values
in a topological linear space $\vecV$ is said to be Cauchy
if for each neighbourhood $\U$ of $\0 \xeof \vecV$
there is an integer $N_{\U}$
such that $x_n - x_m \in \U$ for all $m,n > N_{\U}$.
Then $\vecV$ is called \textit{sequentially complete},
if every Cauchy sequence converges to a unique element $x \xeof \vecV$.
In the setting of metric spaces,
however,
the definition of Cauchy sequence may be simplified by means of the metric:
$\{x_n\}_{n \in \N}$ is Cauchy,
if for each $\epsilon > 0$ there is an integer $N_{\epsilon}$ 
such that
$x_m - x_n \xeof B_{\epsilon}(\0)$ for all $m,n > N_{\epsilon}$,
where $B_{\epsilon}(\0) \eqdef \{x \xeof \vecV: \d(x,\0) < \epsilon\}$
is the open ball with radius $\epsilon$ and center $\0 \xeof \vecV$.
\end{itemize}

\subparagraph{}
For a topological linear space being metrizable and sequentially complete
at the same time,
both concepts of completeness are compatible to each other,
that is,
such a space is complete as a metric space
\iff
it is sequentially complete as a topological linear space.
This equivalence is founded by the fact
that a metric space fulfills the \textit{first countability axiom}.

\paragraph{} %- continuous embeddings
In the sequel we also use the notion of \textit{continuous embedding}
of a topological linear space $\banX$ into a topological linear space $\banY$,
which is nothing but a continuous and \textit{injective} linear mapping
$\iota\:\banX \xto \banY$\footnote{
\,The linearity of $\iota$ means that
$\iota(\alpha x + \beta y) \eq \alpha \iota(x) + \beta \iota(y)$ holds
for all $\alpha, \beta \xeof \F$ and $x,y \xeof \banX$.}.

\subparagraph{} %- linear homeomorphisms
If $\iota$ is surjective and $\iota^{-1}$ is continuous as well,
then both mappings are called \textit{linear homeomorphisms}
between $\banX$ and $\banY$.
By $\banX \cong \banY$ we mean
that $\banX$ and $\banY$ are linear-homeomorphic to each another.

\section{Seminorms and Locally Convex Spaces}
%------------------------------------------------------------------------------%
The class of \textit{locally convex spaces} can be seen as
the most important one within all kinds of topological linear spaces.
The topology $\topT_{\,\vecV}$ of a locally convex space $\vecV$ distinguishes
by the fact
that its origin $\0$
(and thus any other element in $\vecV$) owns a neighbourhood base
consisting of \textit{convex}, \textit{balanced} and \textit{absorbing} sets.

\subparagraph{} %- seminorms
It turns out
that each such pair $(\vecV,\topT_{\,\vecV})$ gives rise
to a family of \textit{seminorms} on the space $\vecV$.
By the concept of seminorm
we mean a mapping $\abs{\,\cdot\,} \:\vecV \xto \R$ having the properties

\begin{itemize}[leftmargin=55pt,itemsep=2pt]
\item[(SN-1)] \quad
$\abs{\lambda x} \eq \abs{\lambda} \cdot \abs{x} \for \lambda \xeof \F$
\textit{(homogenity)},

\item[(SN-2)] \quad
$\abs{x+y} \, \leq \, \abs{x} + \abs{x} \for x,y \xeof \vecV$
\textit{(subadditivity)},
\end{itemize}

which especially implies $\abs{x} \leq 0$ for $x \xeof \vecV$
and $\abs{\0} \eq 0$.

\subparagraph{} %- Minkowski functional
Moreover,
if $\U \xsubset \vecV$ forms a convex, balanced and absorbing neighbourhood
of the origin $\0$,
then the \textit{Minkowski functional}
\[
\begin{cases}[1.3]
\quad \abs{\,\cdot\,} _{\,\U}\:\vecV \xto \R_0^+,\\
\quad x \mapsto \abs{x} _{\,\U} \, \eqdef \,
\mbox{inf} \, \{\lambda \geq 0: x \in \lambda \, \U\} \\
\end{cases}
\]
fulfills all the properties of a seminorm.
On the other hand,
a family $\famF$ of seminorms defined on a linear space $\vecV$
leads in a natural way to a locally convex topology $\topT_{\,\vecV}$:
it is the coarsest topology on $\vecV$
such that all seminorms in $\famF$ remain continuous.

\subparagraph{}
In summary,
there is a one-to-one correspondence
between linear spaces endowed with locally convex topologies
and those endowed with families of seminorms.
Notice that a locally convex topology is Hausdorff
iff
the family of seminorms is \textit{separating},
{\ie} for each element $x \neq 0$
there is at least one seminorm $\abs{\,\cdot\,} _{\alpha}$
such that $\abs{x} _{\alpha} \neq 0$.

\paragraph{} %- F-spaces and Frechet spaces
There are a two subclasses of locally convex spaces,
we want to review in the following:

\begin{itemize}[leftmargin=12pt,itemsep=5pt]
\item[1.]
A topological linear space $\vecV$ is called \textit{F-space},
if it is also a \textit{complete} metric space with translation-invariant
metric $\d$,
{\ie}
\[
\d(x,y) \= \d(x+z,y+z) \; \fa x,y,z \in \vecV,
\]
and the metric topology coincides with the given topology on $\vecV$.

\item[2.]
If the topology of the F-space $\vecV$ is locally convex,
then $\vecV$ is said to be a \textit{{\FRECHET} space}\footnote{
A metrizable locally convex space that is not necessarily complete,
is called a \textit{pre-{\FRECHET} space}.}.
Alternatively,
a locally convex space is {\FRECHET}
\iff
it is complete and its family of seminorms is mostly countable.
In this case
we may define a translation-invariant metric on a {\FRECHET} space $\vecV$
according to
\[
\d(x,y) \, \eqdef \,
\sum_{n=1}^{\infty} p_n \cdot \frac{\abs{x-y} _n}{1+ \abs{x-y} _n}
\quad (x,y \in \topV)
\]
where $\{p_n\}_{n=1}^{\infty}$ is any sequence of positive real numbers,
whose infinite sum converges.
\end{itemize}

\section{Normed Linear Spaces}
%------------------------------------------------------------------------------%
Based on the preceded results,
a normed linear space is nothing but a pre-{\FRECHET} space,
whose family of separating seminorms consists of only one element
$\norm{\,\cdot\,} \eqdef \abs{\,\cdot\,} _0$,
which of course must be a separating one,
that is,
$\norm{x} \eq 0$ implies $x \eq \0$.

\subparagraph{}
It is possible to introduce the concepts of norm and normed linear space
in a direct way, 
{\ie} without using the theory on families of seminorms
and locally convex spaces.
Following this approach,
a \textit{normed space} $\banX$ is a $\F$-linear space
that differs from other abstract linear spaces
by the presence of a particular mapping $\norm{\,\cdot\,}: \banX \xto \R$,
denoted the \textit{norm} on $\banX$,
which fulfills for all $x,y \xeof \banX$ and $\lambda \xeof \F$

\begin{itemize}[leftmargin=55pt,itemsep=2pt]
\item[(N-1)] \quad
$\norm{\lambda \, x} \= \abs{\lambda} \cdot \norm{x}$,

\item[(N-2)] \quad
$\norm{x+y} \, \leq \, \norm{x} + \norm{y}$
\; \textit{(triangle inequality)},

\item[(N-3)] \quad
$\norm{x} \geq 0$ and $\norm{x} \eq 0 \Leftrightarrow x \eq \0$.
\end{itemize}

\paragraph{}
Linear spaces can be furnished with different norms.
Then two arbitrary norms $\norm{\,\cdot\,}_{\alpha}$
and $\norm{\,\cdot\,}_{\beta}$ defined on a linear space $\banX$
are called \textit{equivalent} to each another,
if there are constants $c_1,c_2 > 0$ such that
$
c_1 \norm{x}_{\alpha} \leq \norm{x}_{\beta}  \leq c_2 \norm{x}_{\alpha} 
$
for all $x \xeof \banX$.

\subparagraph{} %- norm topologies
Each normed linear space $\banX$ can be equipped
with a very natural topology $\topT$,
called the \textit{norm topology} on $\banX$.
It can be defined either by the assertion
that $\topT \eqdef \topT_{\norm{\cdot}}$ is the smallest topology on $\banX$
such that the norm mapping remains continuous,
or it is defined to be the metric topology $\topT_d$,
when $\banX$ is considered as a metric space with metric function
$\d(x,y) \eqdef \norm{x-y}$ for all $x,y \xeof \banX$.
In fact,
both definitions lead to the same collection of open subsets in $\banX$,
so $\topT \eqdef \topT_{\norm{\cdot}} \eq \topT_d$,
which justifies to speak of \textit{the} topology of a normed linear space.
Clearly,
equivalent norms on $\banX$ own one and the same topology.
For example,
if $\banX$ is a finite-dimensional linear space,
then all norms on $\banX$ are equivalent and each of them
generates the so-called \textit{standard topology} on $\banX$.

\paragraph{} %- Riesz's lemma
The following assertion,
known as \textit{Riesz's lemma},
quotes
that if $\banX$ is a normed space
and $U \subsetneq \banX$ is a proper closed subspace,
then for every $0 < \alpha < 1$ there is a unit vector $v \xeof \banX$
such that $\norm{v{-}u} \geq \alpha$ for all $u \xeof U$.
This is a useful tool,
which can help to prove a couple of other theorems in functional analysis.

\subparagraph{} %- norm-isomorphisms
if $\banX$ and $\banY$ are normed linear spaces
and $\iota$ is a \textit{norm-preserving} linear homeomorphism
between $\banX$ and $\banY$,
that is,
$\norm{\iota x}_{\banY} \eq \norm{x}_{\banX}$ for all $x \xeof \banX$,
we denote this mapping a \textit{norm-isomorphism} between $\banX$ and $\banY$,
and in this case both spaces are called \textit{norm-isomorphic} to each other
(we again use the notion $\banX \cong \banY$ in this case).

\subparagraph{} %- separable spaces
A normed linear space $\banX$ is called \textit{separable}
if there is a mostly countable subset lying dense in $\banX$.

\paragraph{} %- norm convergence
We finally introduce the concept of convergence
for sequences in a \textit{normed} linear space $\banX$.
Any sequence $\{x_n\}_{n \in \N}$ in $\banX$
is called \textit{strongly convergent}
or simply \textit{norm convergent} to $x \xeof \banX$,
if it converges to $x$ {\wrt} the norm topology,
that is,
if $\lim_{\,n \xto \infty} \norm{x_n-x} \eq 0$.

\section{Strictly and Uniformly Convex Spaces}
%------------------------------------------------------------------------------%
In this paragraph we consider certain normed linear spaces,
in which the best approximation of arbitrary elements
by other elements contained in a convex subset is always possible.

\paragraph{} %- uniformly convex spaces
A \textit{uniformly convex space} is a \textit{real} normed linear space $\banX$
having the following property:
for every $\epsilon > 0$ there is $\delta > 0$ 
such that
for any elements $x,y$ of the closed unit ball $B_1(\0)$ of $\banX$
the validness of the inequality $\norm{(x+y)/2} > 1- \delta$
implies $\norm{x-y} < \epsilon$.
In other words,
if the center $(x+y)/2$ of the line segment $[x,y]$
between elements $x,y$ belonging to the \textit{unit sphere}
$S_1(\0) \eqdef \partial B_1(\0)$ tends to $S_1(\0)$,
then the $x$ and $y$ tend together.
This property is possibly lost
when switching to another but equivalent norm on $\banX$.

\subparagraph{} %- strictly convex spaces
A normed linear space $\banX$ is called \textit{strictly convex},
if the conditions
$\norm{x} \leq 1$, $\norm{y} \leq 1$ and $x \neq y$
imply
$\norm{x+y}/2< 1$.
This means,
that the \textit{inner} line segment $(x,y)$
between $x,y \xeof S_1(\0)$ lies completely in the interior of $B_1(\0)$,
that is
$\norm{\,\alpha x + (1-\alpha)\,y\,} < 1$ for $0 < \alpha < 1$.
Each linear subspace of a uniformly convex Banach space owns this feature.
In the light of the forthcoming theorem \ref{I_APPROXIMATION_THEOREM},
it is an important fact
that uniformly convex spaces are also strictly convex.

\paragraph{} %- approximation theorem
Recall that a subset $\convC$ of a linear space is called \textit{convex},
if the line segment $[x,y]$ of any pair of points $x,y \xeof \convC$
belongs to $\convC$ again.
Then many approximation problems may be formulated as follows:
it is to find an element $y_0$ in a convex set $\convC \subset \banX$
such that $y_0$ has minimal distance to a given element $y \notin \convC$.
The following statement,
known as the \textit{approximation theorem},
provides sufficient conditions regarding existence and uniqueness
of a solution for that problem:

\begin{theorem1}\label{I_APPROXIMATION_THEOREM}
Let $(\banX,\norm{\,\cdot\,})$ be a real normed linear space,
let $y \xeof \banX$ and
$\convC \subset \banX$ be a \textit{closed} and \textit{convex} subset.
Then we have:

\begin{itemize}[leftmargin=17pt,itemsep=5pt]
\item[(a)] \textit{Uniqueness}:
if $\banX$ is \textit{strictly convex},
then there is mostly one element $y_0 \xeof \convC$
such that $\norm{y_0-y} \eq \inf_{x \xeof \convC} \norm{y-x}$.

\item[(b)] \textit{Existence}:
if $\banX$ is a \textit{uniformly convex Banach space}
then there is a unique element $y_0 \xeof \vecV$
such that $\norm{y_0-y} \eq \inf_{x \xeof \convC} \norm{y-x}$.
\end{itemize}
\end{theorem1}

\section{Banach and Hilbert Spaces}
%------------------------------------------------------------------------------%
A normed linear space $\banX$ is said to be \textit{sequentially complete}
(or briefly \textit{complete}),
if for any Cauchy sequence $\{x_n\}_{n \in \N}$ in $\banX$
there is $x \xeof \banX$
such that $\lim_{\,n \xto \infty} \norm{x_n-x} \eq 0$.
Equivalently,
the space $\banX$ is complete
\iff
for every absolutely convergent series
$\sum_{n=1}^{\infty} \norm{x_n} < \infty$ there is $x \xeof \banX$
such that $\norm{x-x_n} \xto 0$ for $n \xto \infty$.
Each complete normed linear space is denoted a \textit{Banach space}.
Like in the non-complete case,
a Banach space $\banX$ with norm $\norm{\,\cdot\,}$
is nothing but a {\FRECHET} space,
whose separating family of seminorms contains only this norm.

A complex (real) \textit{Hilbert space} differs
from an arbitrary Banach space by the presence of an \textit{inner product},
which is nothing but a symmetric and positive sesquilinear (bilinear) form
on such a space.
This form defines the norm and the metric topology of the Hilbert space,
but also determines the inner geometry to be close
to that of an Euclidean space.

\subparagraph{} %-- inner geometry of Hilbert spaces
Let us make these assertions more precise.
Given a real or complex linear space $\hilH$,
an \textit{inner product} in $\hilH$ is defined to be a mapping
$\scalar{\,\cdot\,}{\,\cdot\,}\: \hilH \xtimes \hilH \xto \F$,
which is

\begin{itemize}[leftmargin=35pt]
\item[(I-1)]
linear in the first argument,

\item[(I-2)]
either linear or antilinear in the second argument (conditional upon,
whether $\F$ is the real or the complex number field),
\end{itemize}

together with the additional properties

\begin{itemize}[leftmargin=35pt]
\item[(I-3)]
$\scalar{y}{x} \eq           \scalar{x}{y}$  ($\F \eq \R$) rsp.
$\scalar{y}{x} \eq \overline{\scalar{x}{y}}$ ($\F \eq \C$),
\item[(I-4)]
$\scalar{x}{x} \geq 0$ for all $x \neq 0$,
\item[(I-5)]
$\scalar{x}{x} \eq 0$ implies $x \eq 0$.
\end{itemize}

\subparagraph{}
Property (I-3) is denoted the \textit{Hermitean symmetry}
of $\scalar{\,\cdot\,}{\,\cdot\,}$,
while (I-4) and (I-5) together mean \textit{positive definiteness}.
The presence of the inner product gives rise to define a norm in $\hilH$
according to
\begin{equation}\label{I_DERIVED_NORM}
\norm{x} \, \eqdef \, \scalar{x}{x} ^{1/2} \quad \quad (x \in \hilH),
\end{equation}
so each linear space furnished with an inner product is normed
in a natural manner.
The inner product and the norm are coupled by
the \textit{Cauchy-Schwarz} \textit{inequality}
\[
\abs{\scalar{x}{y}} \, \leq \,
\norm{x} \cdot \norm{y} \quad \quad (x,y \in \hilH),
\]
which implies that the mapping
$\scalar{\,\cdot\,}{\,\cdot\,}\:\hilH \xtimes \hilH \xto \F$ is continuous.
An inner product space $\hilH$ is called a \textit{Hilbert space}
or shortly \textit{Hilbertian},
if it is complete {\wrt} the norm defined by \eqref{I_DERIVED_NORM}.

\subparagraph{} %- parallelogram equation and polarization formula
The pleasant geometric properties of inner product spaces
may be attributed to the validness of the \textit{parallelogram equation}
\begin{equation}\label{I_PARALLELOGRAM_EQUATION}
2 \norm{x}^2 + 2 \norm{y}^2 \= \norm{x+y}^2 + \norm{x-y}^2
\qfa x,y \xeof \hilH.
\end{equation}
In general,
the norm $\norm{\,\cdot\,}$ on a linear space $\banX$
can be induced by an inner product
\iff \eqref{I_PARALLELOGRAM_EQUATION} holds.
In this case and if $\banX$ is a complex linear space,
the inner product can be recovered from the norm
by the \textit{polarization formula}
\[
\scalar{x}{y}
\=
\frac{1}{4} \,
(\,
\norm{x + y}^2 - \norm{x-y}^2 +i \norm{x+iy}^2 - i \norm{x-iy}^2
\,).
\]
If $\banX$ is a real space,
however,
the polarization formula simplifies to
\[
\scalar{x}{y} \=
\frac{1}{4} \,
(\,
\norm{x + y}^2 - \norm{x-y}^2
\,).
\]

\paragraph{} %- properties of Hilbert spaces
In the following we want to point out
that Hilbert spaces differ from general Banach spaces
by a couple of desirable properties:

\begin{itemize}[leftmargin=12pt,itemsep=5pt]
\item[1.] \textit{Orthogonality}.
Two elements $x,y \xeof \hilH$ are called \textit{orthogonal} to each other,
if $\scalar{x}{y} \eq 0$ (notation: $x \bot y$).
Any subset $\orthO \xsubset \hilH$ is said to be \textit{orthogonal}
or an \textit{orthogonal system},
if all elements in $\orthO$ are orthogonal to each other.
Finally,
the \textit{orthogonal space} $\setA^{\bot}$ of a subset $\setA \xsubset \hilH$
contains exactly the elements $y \xeof \hilH$,
which are orthogonal to all $x \xeof \setA$,
and forms a closed linear subspace of $\hilH$.

\item[2.] \textit{Projection Theorem}.
Let $\setA \xsubset \hilH$ be a \textit{closed} linear subspace.
Then for each $x \xeof \hilH$ there is exactly one $x_0 \xeof \setA$
such that
\[
\norm{ x - x_0 } \= \inf \, \{ \, \norm {x-y}: y \in \setA \,\}.
\]
The element $x_0$ is called the \textit{projection} of $x$ onto $\setA$
and minimizes the distance between $\setA$ and $x$.
The associated linear mapping
$\pi_{\setA} \: x \xeof \hilH \mapsto x_0 \xeof \setA$
is also denoted the \textit{projection} of $\hilH$ onto $\setA$.
In general,
by a \textit{projection} is meant
any linear mapping $\pi\:\hilH \xto \hilH$ satisfying
$\pi \circ \pi \eq \pi$ and
$\scalar{\pi x}{y} \eq \scalar{x}{\pi y}$ for all $x,y \xeof \hilH$.

\item[3.] \textit{Decomposition Theorem}.
Using the notation of orthogonal space $\setA^{\bot}$,
the projection of an element $x$ onto $\setA$ is characterized as follows:
$x_0 \eq \pi_{\setA} x$ iff $x-x_0 \xeof \setA^{\bot}$.
Hence each $x \xeof \hilH$ can be uniquely represented by the formula
$x \eq x_0 + x^{\bot}$,
where $x_0 \xeof \setA$ and $x^{\bot} \xeof \setA^{\bot}$,
so the entire space $\hilH$ is isometrically isomorphic
to the direct sum $\setA \, \oplus \, \setA^{\bot}$
of the closed subspaces $\setA$ and $\setA^{\bot}$.
This fact is known as \textit{decomposition theorem} in Hilbert spaces.
In other words,
each closed linear subspace $\setA$ of $\hilH$
is \textit{topologically complementable},
that is,
to each closed subspace $\setA$
there is a second closed subspace $\setB \xsubset \hilH$
such that $\setA \cap \setB \eq \{\0\}$ and $[\setA \cup \setB] \eq \hilH$.
Of course,
a closed subspace $\A$ is topologically complemented
by its orthogonal space. %$\setA^{\bot}$.

\item[4.] \textit{Existence of Orthonormal Bases}.
An orthogonal system $\orthO \xsubset \hilH$ is denoted
an \textit{orthonormal system},
if $\norm{x} \eq 1$ for all $x \xeof \orthO$.
However,
if $\orthO$ contains only \textit{finitely} many elements,
then the span $[\orthO]$ forms a \textit{closed} subspace in $\hilH$,
and the projection $\pi_{[\orthO]}$ of any element $x$ onto $[\orthO]$
can be expressed by the formula
\[
\pi_{[\orthO]} \, x \= \sum_{i=1}^n \, \scalar{x}{e_i} \, e_i.
\]
However,
if $\orthO$ consists of arbitrarily many elements,
the inequality $\scalar{x}{e_{\alpha}} \neq 0$ can only be valid
for mostly countable many elements $e_{\alpha} \xeof \orthO$,
which in conclusion enables to add the coefficients $\scalar{x}{e_{\alpha}}$
together about all elements of $\orthO$.

\subparagraph{} %- orthonormal bases
An orthonormal system $\orthO$ is called
an \textit{orthonormal base} in $\hilH$,
if its span $[ \orthO ]$ lies dense in $\hilH$.
In this case we may represent each $x \xeof \hilH$ according to
\begin{equation}\label{I_FOURIER_COEFFICIENTS}
x \= \sum_{e_{\alpha} \in \, \orthO} \, \scalar{x}{e_{\alpha}} \, e_{\alpha},
\end{equation}
which is primarily a formal sum,
but consists of mostly countable many summands.
The inner products $\scalar{x}{e_{\alpha}}$ are called
the \textit{Fourier coefficients} of $x$ {\wrt} $\orthO$.
The right hand side of \eqref{I_FOURIER_COEFFICIENTS}
consists of \textit{mostly countable many} non-zero summands,
and this series is invariant
by realignments of the elements $e_{\alpha} \xeof \orthO$.

\subparagraph{}
There is the need to determine,
whether a given orthonormal system $\orthO$ forms an orthonormal base.
In fact,
it can be shown that
the following statements are equivalent:

\begin{itemize}[leftmargin=20pt]
\item[(a)]
$\orthO$ is an orthonormal base in $\hilH$.

\item[(b)]
$\orthO$ is \textit{complete},
that is,
$\scalar{x}{e} \eq 0$ for all $e \xeof \orthO$ implies $x \eq \0$.

\item[(c)]
$\sum_{e_{\alpha} \in \,\orthO} \;
\abs{ \scalar{x}{e_{\alpha}} } ^2 \eq \norm{x}^2$ for all $x \xeof \hilH$
(\textit{Parseval}\textit{'s equation}).
\end{itemize}

\paragraph{} %- existence of orthonormal bases
It is also nearby to ask for the existence of orthonormal bases.
First of all,
the collection of all orthonormal systems in a Hilbert space $\hilH$
forms a partial ordered set,
and under the assumption that \textsc{Zorn}{'s lemma}\,\footnote{\,
\,The \textit{axion of choice} claims that
if $\cal{F}$ is a disjoint family $F_{\alpha}$ ($\alpha \in \J$)
of non-empty sets,
then there is a set $A$ that shares exactly one element
with each other set $F_{\alpha} \in \cal{F}$.}
is valid,
it contains a \textit{maximal} element $\orthO$
that cannot be increased,
hence each Hilbert space $\hilH \neq \{\0\}$ owns an orthonormal base.

\subparagraph{} %- Hilbert space dimension
Furthermore,
if $\orthO_1$ and $\orthO_2$ are orthonormal bases in $\hilH$,
then there is a bijective mapping between them,
thus they own identical cardinal numbers.
This legitimates us to introduce the notation
of \textit{Hilbert space dimension} $\abs{\hilH} $
being the cardinal number of any orthonormal base in $\hilH$.

\subparagraph{} %- separable Hilbert spaces
The space $\hilH$ is separable \iff $\abs{ \hilH} \leq \abs{\N} $.
Indeed,
each orthonormal base in $\hilH$ forms a mostly countable set in this case,
and $\hilH$ proves to be norm-isomorphic to the Hilbert space $\ell^2$
of square summarizable sequences $f\:\N \xto \F$.
\end{itemize}

\section{Spaces of Scalar-Valued Functions and Measures}
%------------------------------------------------------------------------------%
This paragraph forms an enumeration of well-known and important
topological linear spaces of functions or measures.
It is assumed that the involved functions are either real or complex-valued.

\begin{itemize}[leftmargin=12pt,itemsep=5pt]
\item[1.] \textit{Spaces of continuous functions.}\\
For any topological space $\topX$
the function spaces $C_b(\topX)$ and $C_0(\topX)$,
both endowed with the \textit{supremum norm}
\[
\norm{f} \eqdef \sup_x \, \abs{f(x)} \,,
\]
where $x$ runs through the elements of the underlying space $\topX$.
Both are complete and hence examples of Banach spaces.
If the space $\topX$ is compact,
then $C(\topX)$,
the space of continuous functions $f\:\topX \xto \C$,
forms a Banach space as well.

\item[2.] \textit{The space $\MS{X}$.}\\
Recall from measure theory
that a \textit{signed measure},
defined on a measurable space $(X,\salgB)$,
is a $\sigma$-\textit{additive} set function
\[
\mu: \salgB \to (-\infty,+ \infty] \quad \mbox{ or } \quad
\mu: \salgB \to [-\infty,+ \infty).
\]
If $\mu$ is a \textit{finite signed measure}
(the values $+\infty$ und $-\infty$ are excluded in this case),
it can be represented as the difference of positive measures $\mu_+$
(\textit{Hahn-Jordan decomposition theorem}).
In this case
the positive measure $\abs{\mu} := \mu_+ + \mu_-$ is called the
\textit{variation}\ of $\mu$,
and the mapping $\norm{\mu} := \abs{\mu} (X)$ is said to be
the \textit{total variation} of $\mu$.
It can be shown 
that the total variation constitutes a norm on the space $\MS{X}$
of finite signed measures,
which is even complete and hence forms a Banach space.

\item[3.] \textit{The space $\MC{X}$.}\\
Similarly to finite signed measures,
a \textit{complex measure} on a measurable space $(X,\salgB)$
is a $\sigma$-additive set function $\mu:\salgB \to \C$.
The set $\MC{X}$ of complex measures on $(X,\salgB)$
actually forms a $\C$-linear space.
It is possible to attach to each complex measure $\mu$
an arbitrary positive measure $\abs{\mu} $,
called the \textit{total variation} of $\mu$.
It is the mapping
$\abs{\mu} :\salgB \to [0,\infty]$,
which assigns to each measurable set $A \xeof \salgB$ the value
\[
\abs{\mu} (A)
\, := \,
\sup \, \sum_{n=1}^{\infty} \, \norm{\,\mu(A_n)},
\]
where $(A_n)$ may be any mostly countable and disjoint decomposition
of $A$ into measurable subsets of $A$.
Indeed,
the mapping $\abs{\mu} $ is a positive measure,
and a complex measure $\mu$ is called of \textit{bounded variation},
if $\abs{\mu} (X) < \infty$.
The mapping 
\[
\mu \mapsto \norm{\mu}
\, := \,
\abs{\mu} (X)
\]
satisfies all conditions of a norm,
hence the pair $(\MC{X} ,\norm{\cdot})$ forms a normed linear space.
Like the space $\MS{X}$ of finite signed measures,
the space $\MC{X}$ is complete and thus a Banach space.

\item[4.] \textit{The spaces $BV[a,b]$ and $AC[a,b]$.}\\
Let $BV[a,b]$ be the linear space consisting
of all functions of bounded variation.
This space can be equipped with the norm
\[
\norm{f}_{bv} \, \eqdef \,
\abs{ f(a) } \, + \,
\sup_{P} \sum_{i=0}^{n_P-1} \abs{f(x_{i+1}-f(x_i)} ,
\]
where the supremum is taken over all finite partitions of $[a,b]$.
The space $BV(a,b)$ forms a Banach space in this norm.

\subparagraph{} %- the space AC[a,b]
A measurable function $f\:[a,b] \xto \R$ is called 
\textit{absolutely continuous},
if for each $\epsilon > 0$ there is $\delta > 0$ 
such that for a finite decomposition
\[
x_0=a < x_1 < \ldots < x_n=b \mbox{ mit }
\sum_{k=1}^n \abs{x_k-x_{k-1}} < \delta
\]
of the interval $[a,b]$
the condition $\sum_{k=1}^n \abs{f(x_k)-f(x_{k-1})} < \delta$ is valid.
Every absolutely continuous function is of bounded variation.
The etenge $AC[a,b]$ of absolutely continuous functions on $[a,b]$
endowed with the norm $\norm{\,\cdot\,}_{BV}$
forms a closed linear subspace of $BV[a,b]$ and hence 
is a Banach space in its own right. 
\end{itemize}

\paragraph{} %- spaces of functions defined on a domain G
Let us continue with the consideration of function spaces,
where the underling set $\G$ is a \textit{domain} in $\Rd$,
that is,
an open and connected set within an Euclidean space $\Rd$,
which may be bounded, unbounded or even the space $\Rd$ itself.

\subparagraph{} %- multi-indices
We will also work in the sequel with \textit{multi-indices},
which are nothing but $d$-tuples
$\alpha \eq (\alpha_1,\ldots,\alpha_d)$
with values $\alpha_i \xeof \N_0$
and modulus $\abs{\alpha} \eqdef \alpha_1 + \ldots + \alpha_d$.
They are mostly used to shorten multiple partial derivatives
$
\PA f \eqdef \partial_1^{\alpha_1}
         \partial_2^{\alpha_2} \ldots
         \partial_d^{\alpha_d} f
$
of a sufficiently often differentiable and complex-valued function $f$,
but also to denote the multiple powers
$\xi^{\alpha} \eqdef \xi_1^{\alpha_1}
\cdot \xi_2^{\alpha_2} \cdot \ldots \cdot \xi_d^{\alpha_d}$
of a vector $\,\xi \eq (\xi_1,\xi_2,\ldots,\xi_d) \xeof \C^d$.

\subparagraph{}
\begin{itemize}[leftmargin=12pt,itemsep=5pt]
\item[5.] \textit{Spaces of continuously differentiable functions}.\\
The linear subspace $C^k(\clG) \subset C(\clG)$
consists of those continuously differentiable functions $f\:\G \xto \F$,
whose partial derivatives $\partial^{\alpha} f$ are bounded on $\clG$
for every multi-index $\alpha$ with $\abs{\alpha} \leq k$.
It is complete if endowed with the norm
\[
\norm{f}_k
\, \eqdef \,
\sum_{\abs{\alpha} \leq k} \,
\sup_{x \in \clG } \; \abs{\partial^{\alpha} f(x)} .
\]

\item[6.] \textit{{\HOELDER} spaces}.\\
It has been shown
that the spaces $C^k(\clG)$ are not very convenient
for the investigation of partial differential equations,
but the so-called \textit{{\HOELDER} spaces} $C^{k,\gamma}(\clG)$ are so.
Preliminary,
a function $f\:\G \xto \F$ is called \textit{{\HOELDER}-continuous}
{\wrt} the parameter $\gamma \xeof (0,1]$,
if its \textit{{\HOELDER}-norm} is finite,
that is,
\[
\abs{f} _{\,\gamma} \, \eqdef \,
\sup_{x,y \, \in \, \G, \, x \neq y} \,
\frac{\abs{f(x) - f(y)} } {\abs{x-y} ^{\gamma}} \, < \, \infty.
\]
Let $C^{k,\gamma}(\clG)$ be the linear space
of functions $f \xeof C^k(\clG)$,
whose partial derivatives $\PA f$
are {\HOELDER}-continuous for all $\abs{\alpha} \eq k$.
Then the \textit{{\HOELDER} norm} $\abs{f} _{k,\gamma}$
of a function $f \xeof C^{k,\gamma}(\clG)$ is defined
by the sum of the norm $\norm{f}_k$
and the {\HOELDER} norms of the partial derivatives $\PA f$
for $\abs{\alpha} \eq k$:
\[
\abs{f} _{k,\gamma}
\, \eqdef \,
\norm{f}_k
\, + \,
\sum_{\abs{\alpha} \eq k} \abs{\,\PA f} _\gamma.
\]
The space $C^{k,\gamma}(\clG)$ is complete {\wrt} this norm.
Notice that for $k \eq 0$ and $\gamma \eq 1$ we obtain the space
$C^{0,1}(\clG)$ of \textit{Lipschitz-continuous} functions.

\subparagraph{}
If $\G$ is a bounded domain, we have for $0 < \alpha < \beta \leq 1$
\[
C^1(\clG) \subsetneq
C^{0,\beta}(\clG) \subsetneq
C^{0,\alpha}(\clG) \subsetneq
C(\clG).
\]

\item[7.] \textit{Spaces of smooth functions}.\\
Next we consider three function spaces
consisting of indefinitely differentiable functions $f\:\G \xto \F$
that frequently occure in applied analysis
and constitute {\FRECHET} spaces respectively
complete locally convex linear spaces.

\begin{itemize}[leftmargin=17pt]
\item[(a)]
The space $\SF$ of functions $f\:\G \xto \F$ being indefinitely differentiable.
Unfortunately,
there is no norm that makes $\SF$ into a complete normed space,
but after all it is possible to define a countable family $\famF$
of seminorms to make $\SF$ into a {\FRECHET} space.
For example,
$\famF$ can be given by the mappings
\[
p_{N,n} \, \eqdef \,
\max_{x \in \compK_n} \;
\{ \, \abs{\, (\PA \phi)(x) \,} : \abs{\alpha} \leq N \}
\quad \quad (N,n \in \N)\,,
\]
where $(\compK_n)$ is a sequence of compact subsets of $\G$
such that
\[
\compK_1 \subset \compK_2 \ldots \subset \G \, \mbox{ for } \, n \in \N,
\quad
\bigcup_{n \in \N} \, \compK_n \= \G,
\]
and each $\compK_n$ lies in the interior of $\compK_{n+1}$
(such a sequence is called a \textit{compact exhaustion} of $\G$).
The space $\SF$ is often denoted $\E(\G)$,
when considered in the locally convex topology induced by the totality
of seminorms $p_{N,n}$.
It is complete {\wrt} this topology,
so it forms a {\FRECHET} space.

\item[(b)]
The subspace $\TF \subset \SF$ of smooth functions
having \textit{compact} support
\[
\supp{\,\phi} \, \eqdef \, \overline{\{x \in \G: \phi(x) \neq 0\}},
\]
usually denoted the \textit{space of test functions}.
This function space plays an outstanding role in analysis,
since it lies dense in a multiplicity of other function spaces.
Unfortunately,
there is neither a countable family of seminorms
making $\TF$ into a {\FRECHET} space nor is $\TF$ completely metrizable.

\subparagraph{}
But it is still possible to define a locally convex topology $\topT$ on $\TF$
such that $\TF$ is complete in the sense of a topological linear space.
Instead of saying how the open sets of this topology $\topT$ look like,
let us define $\topT$ by means of a suitable notion of convergence,
known as $\D$-convergence in $\TF$:\,
a sequence $(\phi_n)$ in $\TF$ is convergent to $\phi \xeof \TF$,
if there is a compact set $\compK \subset \G$
such that
(i)  $\supp{\phi_n} \subset \compK$ for all $n \xeof \N$ and
(ii) for each multi-index $\alpha$ the sequence $(\PA \phi_n)$
     converges to $\PA \phi$ \textit{uniformly} on $\compK$.

The space $\TF$ equipped with the topology
induced by $\D$-convergence is usually denoted $\D(\G)$.
Notice
that we have $\D(\G) \hookrightarrow \E(\G)$,
that is,
the space $\D(\G)$ is continuously embedded into the space $\E(\G)$
{\wrt} the locally convex topologies of these spaces.

\subparagraph{}
We have $\D(\G) \hookrightarrow \E(\G)$,
that is,
the space $\D(\G)$ is continuously embedded into the space $\E(\G)$
{\wrt} their locally convex topologies.

\item[(c)]
The \textit{Schwartz space} $\SRd$,
which is the space of all smooth functions $\phi \xeof C^{\infty}(\Rd)$
having the additional property
\[
\abs{ \PA \phi(x) \cdot x^{\beta}} \xto 0 \,
\for \norm{x}_d \xto \infty
\]
and arbitrary multi-indices $\alpha,\beta$,
where $\norm{x}_d$ is the Euclidean norm of $x \xeof \Rd$.
The elements of $\SRd$ are called \textit{rapidly decreasing}
or briefly \textit{Schwartz functions}.
Hence a Schwartz function and all its derivatives vanish at infinity
faster then the reciprocal of any polynomial.
The most famous example of a Schwartz function is the Hermite function
\[
\cal{H}_0(x) \, \eqdef \, e^{- \frac{\norm{x}^2}{2}} \quad \quad (x \in \Rd).
\]

\subparagraph{}
Rather similar to the space $\SF$,
it is possible to find for the Schwartz space $\SRd$
a family of seminorms as follows:
define for $\phi \xeof C^{\infty}(\Rd)$ and multi-indices $\alpha, \beta$
\begin{equation}\label{I_SRN_FAMILY_OF_SEMINORMS}
p_{\alpha,\beta}(\phi) \, \eqdef \,
\sup_{x \in \Rd} \abs{ \, x^{\alpha} (D^{\beta} \phi) (x) \, } .
\end{equation}

\subparagraph{}
Then $\phi$ is a Schwartz function
iff $p_{\alpha,\beta}(\phi) < \infty$
holds for all indices $\alpha, \beta$.
The collection of mappings $p_{\alpha,\beta}\:\SRd \xto \R$
defines a countable family of seminorms,
which gives rise to a locally convex topology $\topT_{\mathscr{S}}$ on $\SRd$.
Moreover,
since the Schwartz space is complete {\wrt} $\topT_{\mathscr{S}}$,
it actually forms a \textit{{\FRECHET} space}.

\subparagraph{}
Comparing the Schwartz space $\SRd$ with the function spaces
$\D(\Rd)$ and $\E(\Rd)$,
it is easy to see that there is the chain of inclusions
\[
\D(\Rd) \, \hookrightarrow \, \SRd \, \hookrightarrow \, \cal{E}(\Rd).
\]
Moreover,
by comparing the locally convex topologies of these spaces,
it turns out that these inclusions are continuous.
\end{itemize}

\item[8.] \textit{Lebesgue spaces}.\\
To introduce the spaces $L^p(\G)$ for $p \xeof [1,\infty)$,
let $\cal{L}^p(\G)$ be the $\F$-linear space of measurable functions
$f\:\G \xto \F$
such that $\abs{f} ^p$ is Lebesgue-integrable.

\subparagraph{}
We define for $f,g \xeof \cal{L}^p(\G)$ the equivalence relation
$f \sim g$ according to $f \eq g$ almost everywhere.
Then the elements of the space $L^p(\G)$
are the equivalence classes in $\cal{L}^p(\G)$ {\wrt} $\sim$,
that is,
$L^p(\G) \eqdef \cal{L}^p(\G) / \sim$.
The \textit{Lebesgue norm} in $\Lp$ is given by
\[
\norm{f}_p
\, \eqdef \,
\big( \int_{\,\G} \abs{f} ^p \, d\mu \,\big) \, ^{\frac{1}{p}}.
\]
Finally,
the \textsc{Fischer-Riesz} \textit{theorem} claims
that the spaces $\Lp$ are complete {\wrt} the Lebesgue norm.
It turns out that $\Lp$ is separable \iff $1 \leq p < \infty$.

\subparagraph{} %- Lebesgue space L2
The case $p \eq 2$ gives rise to a special attention:
here it is possible to derive the norm $\norm{\,\cdot\,}_2$
from the inner product
\[
\scalar{f}{g} \, \eqdef \, \int_{\,\G} f \, \overline{g} \; d\mu,
\]
thus $\norm{f}_2 \eq \scalar{f}{f} ^{1/2}$
and $\Lh$ becomes into a Hilbert space.

\paragraph{}
For $p \eq \infty$ and any measurable function $f\:\G \xto \F$
we define the mapping
\begin{equation}\label{I_SUPREMUM_NORM}
\norm{f}_{\infty} \, \eqdef \, \inf
\{\,a \geq 0: \abs{f} \leq a \,\, \mu \mbox{-almost everywhere}\,\},
\end{equation}
taking values in the set $\R^+ \cup \{\infty\}$.
Let $\cal{L}^{\infty}(\G)$ be the linear space
of all measurable functions $f\:\G \xto \F$
with $\norm{f}_{\infty} < \infty$
and $\Li$ be the linear space of all equivalence classes
in $\cal{L}^{\infty}(\G)$,
where $f \sim g$ means $f=g$ $\mu$-almost everywhere.
Members of these spaces are called \textit{essentially bounded}.
Similar to the case $p < \infty$,
the space $\Li$ is complete {\wrt} \eqref{I_SUPREMUM_NORM},
but not separable.

\item[9.] \textit{Sobolev spaces}.\\
A function $f\:\G \xto \C$ is said to belong to the space $\Lc$,
if $f$ is Lebesgue-integrable on each compact subset of $\G$.
This space is useful to define the concept of \textit{weak differentiability}.
A function $f \xeof \Lc$ is said to be \textit{weakly differentiable}
{\wrt} a given multi-index $\alpha$,
if there is $g \xeof \Lc$
such that 
\[
\int_{\,\G} g(x) \, \phi(x) \, dx
=
(-1)^{ \abs{\alpha} } \,
\int_{\,\G} f(x) \, \PA \phi(x) \, dx
\for \phi \in \TF.
\]
We denote $\PA f \eqdef g$ in this case,
and weak differentiability arises as a linear mapping
$\PA\:D_{\alpha} \subset \Lc \xto \Lc$,
where each $D_{\alpha}$ is a linear subspace of $\Lc$.
Notice that
the concept of weak derivative is compliant
with that of the classical derivative:
if $f$ has classical derivative $\PA f$,
then $f$ is also weakly differentiable,
and the weak derivative is equal to $\PA f$.

\begin{itemize}[leftmargin=17pt]
\item[(a)] The next function spaces under consideration are those,
whose elements are $p$-integrable and weakly differentiable up to order $k$
at the same time.
More precisely,
for any integer $k$ and real $p \geq 1$,
the \textit{Sobolev space} $\SobW$ is the set
of functions $f$ in the Lebesgue space $\Lp$,
whose weak derivatives $\PA f$ exist and also belong to $\Lp$
for all multi-indices $\alpha$ with $ \abs{\alpha} \leq k$,
briefly written
\[
\SobW
\, \eqdef \,
\{\,f \in \Lp: \PA f \in \Lp, \; \abs{\alpha} \leq k \, \}.
\]
Sobolev spaces are of course linear spaces,
and especially for $k \eq 0$
the spaces $W^{0,p}(\G)$ are nothing but the Lebesgue spaces $\Lp$ themselves.

\subparagraph{} %- Sobolev norms
For functions $f \xeof \SobW$ let the mapping $f \mapsto \abs{f} _{k,p}$
be defined according to
\begin{equation}\label{I_SOBOLEV_NORM}
\abs{f} _{k,p} \, \eqdef \,
\big(
\sum_{\abs{\alpha } \leq k}
\int_{\,\G} \, \abs{\PA f} ^p \, dx
\big) ^{1/p}.
\end{equation}
It fulfills all the properties of a norm,
thus the space $\SobW$ is actually a normed space,
and the mapping defined by \eqref{I_SOBOLEV_NORM}
is usually denoted a \textit{Sobolev norm}.
With respect to the question of completeness,
it is not hard to show that
the spaces $\SobW$ are complete
and therefore constitute Banach spaces as well.

\subparagraph{} %- alternative definition of Sobolev space
There is a second way to define a Sobolev space.
Consider the subspace
\[
C_{k,p}^{\infty}(\G) \, \eqdef \, \{\, f \in \SF: \abs{f} _{k,p} < \infty \,\},
\]
and since this space is \textit{not} complete
{\wrt} the norm $\abs{\,\cdot\,} _{k,p}$,
we set
\[
\cal{W}^{k,p}(\G) \, \eqdef \, \overline{C_{k,p}^{\infty}(\G)}.
\]
The space $\cal{W}^{k,p}(\G)$ is defined by a purely topological operation,
that is,
by taking the closure of a smaller space.
Thus for elements in this space there is no instrinsic notion of derivative,
and one introduces the notion of \textit{strong derivative} as follows:
given $u \xeof \cal{W}^{k,p}(\G)$
and a multi-index $\alpha$ ($\abs{\alpha} \leq k$),
a further function $\PA u$ also belonging to $\cal{W}^{k,p}(\G)$
is said to be the $\alpha$-strong derivative of $u$,
if there is a sequence $\{\phi_n\} \xeof C_{k,p}^{\infty}(\G)$
converging to $u$ such that
\[
\abs{\PA u - \PA \phi_n)} _{k,p} \xto 0 \quad \quad (n \xto \infty).
\]
Sobolev spaces are reflexive for $1 < p < \infty$,
but not for $p \eq 1$.

\paragraph{} %- Meyers-Serrin theorem
The \textsc{Meyers-Serrin} \textit{theorem} claims
$\cal{W}^{k,p}(\G) \cong \SobW$ in the sense of Banach space isomorphisms.
Moreover,
the notions of strong and weak derivative are equivalent
when identifying functions in $\cal{W}^{k,p}(\G)$ and $\SobW$.

\paragraph{} %- embedding theorems for Sobolev spaces
One of the most important properties of Sobolev spaces is given by the fact
that under certain hypotheses they can continuously be embedded
into further spaces and those of type $C^k(\clG)$.
We define for fixed $1 \leq p < \infty$
the \textit{Sobolev numbers} $\theta_{\,k,d,p}$ by
\[
\theta_{\,k,d,p} \eqdef k - d/p
\quad \quad (k,d \in \N, p \geq 1),
\]
which play a crucial role in \textsc{Sobolev}\textit{'s embedding theorem}:

\begin{theorem1}\label{I_SOBOLEV_EMBEDDING_THEOREM_1}
Let $\G \subset \Rd$ be a domain having Lipschitz-continuous boundary.
If $k,m$ are chosen such that $\theta_{\,k,d,p} > m$,
then there are continuous embeddings
$\SobW \hookrightarrow C^m(\clG)$ and
$\SobW \hookrightarrow C^{m,\gamma}(\clG)$ for $0 < \gamma < 1$.
\end{theorem1}

Hence each class $f \xeof \SobW$ has a representative in these spaces,
supposed the Sobolev number is sufficiently large.

\paragraph{} %- embeddings of Sobolev spaces into other Sobolev spaces
In case of $\theta_{\,k,d,p} \leq 0$,
however,
one cannot embed $\SobW$ into $C(\clG)$.
But it is at least possible to find an embedding of $\SobW$
into some lower-dimensional Lebesgue spaces $L^q(\G)$:

\begin{theorem1}\label{I_SOBOLEV_EMBEDDING_THEOREM_2}
For $\G \subset \Rd$,
$k \xeof \N$ and $p,q \geq 1$ such that
\[
1 < q < \frac{dp}{d-kp},
\]
there are continuous embeddings $\SobW \hookrightarrow L^q(\G)$.
\end{theorem1}

\subparagraph{} %- embeddings for fractional Sobolev spaces
Finally,
there are also embeddings of Sobolev spaces into further Sobolev spaces,
where the parameter $p$ is fixed:

\begin{theorem1}\label{I_SOBOLEV_EMBEDDING_THEOREM_3}
If $\G \subset \Rd$ has Lipschitz-continuous boundary,
$s,s' \xeof \N$ and $1 < s < s' < \infty$,
then $W^{s',p}(\G) \hookrightarrow W^{s,p}(\G)$.
\end{theorem1}

\item[(b)] Let us consider in the family $\SobW$ of Sobolev spaces
the case $p \eq 2$ and denote the space $H^k(\G) \eqdef W^{k,2}(\G)$
for $k \xeof \N_0$,
which is often called a \textit{Hilbert-Sobolev space}.

\subparagraph{} %- inner product and norm of Hilbert-Sobolev spaces
Each space $H^k(\G)$ can be endowed with the inner product
\[
\scalar{f}{g}_k \, \eqdef \,
\sum_{\abs{\alpha} \leq k} \int_{\,\G} \;
\PA f \; \PA \overline{g} \, dx\,,
\quad \quad f,g \xeof H^k(\G)
\]
and with norm
$\abs{f} _k \eq \scalar{f}{f}^{1/2}$ for $f \xeof H^k(\G)$,
so $H^k(\G)$ forms a Hilbert space in its own right.

\item[(c)] The space $\TF$ of test functions can surely be endowed
with the norms $\abs{\cdot} _{k,p}$,
but this leads to uncomplete normed spaces.
Hence it is nearby to introduce the spaces $\WN{k,p}{\G}$ for $k \xeof \N$
by completion of the space $\TF$ {\wrt} the Sobolev norms:
\[
\WN{k,p}{\G} \, \eqdef \, \overline{\TF_{\abs{\cdot} _{k,p}}}.
\]
The elements of these spaces can be identified with functions in $\SobW$.
Furthermore,
the spaces $\WN{k,p}{\G}$ form closed linear subspaces
of the Sobolev spaces $\SobW$.
The spaces $\HN{k}{\G} := W^{k,2}{\G}$ are Hilbert spaces
with inner products and norms as defined for $H^k(\G)$.
\end{itemize}
\end{itemize}

\begin{comment}
\paragraph{}
Further normed linear spaces like Bochner spaces
and Sobolev spaces will be introduced later
after having available the notion of \textit{adjoint space} (section III)
and methods for integration and differentiation of vector-valued functions
(section V).
\end{comment}

\section{Hierarchy of Topological and Normed Linear Spaces}
%------------------------------------------------------------------------------%
\paragraph{} %- hierarchy of spaces
The following figure shows the inheritance tree
of topological and normed linear spaces,
where Banach and Hilbert spaces are positioned at the very end
of this tree.
%For the notion of \textit{reflexive space} we refer to section III.

\begin{center}
\scalebox{0.67}{ % optimum = 0.69

\begin{tikzpicture}[node distance=4.0cm,auto]

\node[align=center,rectangle,draw](TLS) at (0,-1.0)
{\;topological linear space\; \\ $(\topV,\topT_V)$};

\node[align=center,rectangle,draw](LCTLS) at (-5,-3.5)
{locally convex \\ \;topological linear space\; \\ $(\topV,\{ p_{\alpha} \}_{\alpha \in \cal{I}})$};

\node[align=center,rectangle,draw](SCTLS) at (5,-3.5)
{\textbf{sequentially complete}\\ \; \textbf{topological linear space} \; \\ $(V,\topT_V)$};

\node[align=center,rectangle,draw](MS) at (-3,-5.9)
{metric space \\ $(\metrS,\topT_d)$};

\node[align=center,thick,rectangle,draw](CMS) at (1,-5.9)
{\textbf{complete} \\ \; \textbf{metric Space} \; \\ $(\metrS,\topT_d)$};

\node[align=center,thick,rectangle,draw](F-S) at (5,-5.9)
{\; \textbf{F-space}\; \\ $(F,\topT_V=\topT_d)$};

\node[align=center,rectangle,draw](PFS) at (-5,-8.0)
{\;pre-{\FRECHET} space\; \\ $(F,\{p_n\}_{n \in \N})$};

\node[align=center,thick,rectangle,draw](FS) at (5,-8.0)
{\;\textbf{{\FRECHET} space}\; \\ $(F,\{p_n\}_{n \in \N})$};

\node[align=center,rectangle,draw](NS) at (-5,-10.5)
{\;{normed linear space}\; \\$(\banX,\norm{\cdot})$};

\node[align=center,thick,rectangle,draw](BS) at (5,-10.5)
{\;\textbf{Banach space}\; \\$(\banX,\norm{\cdot})$};

\node[align=center,thick,rectangle,draw](RBS) at (5,-12.0)
{\;\textbf{reflexive space}\; \\$(\banX,\norm{\cdot})$};

\node[align=center,thick,rectangle,draw](UCS) at (5,-13.5)
{\;\textbf{uniformly convex space}\; \\$(\banX,\norm{\cdot})$};

\node[align=center,rectangle,draw](IPS) at (-5,-15)
{\;{inner product space}\; \\$(\vecH,\scalar{\cdot}{\cdot})$};

\node[align=center,thick,rectangle,draw](HS) at (5,-15)
{\;\textbf{Hilbert space}\; \\$(\hilH,\scalar{\cdot}{\cdot})$};

\draw[->,>={Stealth}, thick] (MS) to (CMS);
\draw[->,>={Stealth}, thick] (MS) to (PFS);
\draw[->,>={Stealth}, thick] (CMS) to (F-S);
\draw[->,>={Stealth}, thick] (TLS) to (SCTLS);
\draw[->,>={Stealth}, thick] (TLS) to (LCTLS);
\draw[->,>={Stealth}, thick] (SCTLS) to (F-S);
\draw[->,>={Stealth}, thick] (F-S) to (FS);
\draw[->,>={Stealth}, thick] (LCTLS) to (PFS);
\draw[->,>={Stealth}, thick] (PFS) to (NS);
\draw[->,>={Stealth}, thick] (NS) to (IPS);
\draw[->,>={Stealth}, thick] (PFS) to (FS);
\draw[->,>={Stealth}, thick] (NS) to (BS);
\draw[->,>={Stealth}, thick] (IPS) to (HS);
\draw[->,>={Stealth}, thick] (FS) to (BS);
\draw[->,>={Stealth}, thick] (BS) to (RBS);
\draw[->,>={Stealth}, thick] (RBS) to (UCS);
\draw[->,>={Stealth}, thick] (UCS) to (HS);

\node[align=center](I) at (0,0.4) {\textit{metric and topological linear spaces}};
\node[align=center](II) at (   0,-16.4) {\textit{normed linear spaces}};
\draw[dashed,thick]        (-7.8,  0.0) rectangle (8.0,-9.0);
\draw[dashed,thick]        (-7.8, -9.5) rectangle (8.0,-16.0);
\end{tikzpicture}
}

\footnotesize{Figure \thepart.1:
              The hierarchy of metric and topological linear spaces}
\end{center}

% BIBLIOGRAPHY
%------------------------------------------------------------------------------%
\vspace{5mm}
\begin{thebibliography}{99}
\begin{small}

\bibitem{I_DS} N.\,Dunford, J.\,T.\,Schwartz:
\textit{Linear Operators I}.
Wiley \& Sons, 1958.

\bibitem{I_RS1} M.\,Reed, B.\,Simon:
\textit{Modern Methods of Mathematical Physics I: {\FA}},
Academic Press Inc., 1980.

%\bibitem{I_Yo} K.\,Yosida:
%\textit{Functional Analysis}.
%Springer, 1980.

\end{small}
\end{thebibliography}

\newpage
\thispagestyle{empty}
\fancyhead[C]{\small{II. Linear Operators in Banach Spaces}}
%------------------------------------------------------------------------------%
\part{Linear Operators in Banach Spaces}                                       %
%------------------------------------------------------------------------------%
The aim of this second part
is to present some fundamental facts on linear operators
acting between Banach spaces.
Furthermore,
we want to introduce various classes of linear operators
by means of the validness of certain topological properties,
which to the end lead to the concept of
\textit{continuous},
\textit{compact},
\textit{closed} and
\textit{closable} operator.
It is also our intention to point out certain relationships
between these notions and to establish some equivalences.

\section{Fundamentals of Operator Theory}
\section{Continuous Linear Operators}
\section{Compact Linear Operators}
\section{Closed and Closable Linear Operators}
\section{The Closed Graph Theorem}

% BIBLIOGRAPHY
%------------------------------------------------------------------------------%
\vspace{5mm}
\begin{thebibliography}{99}
\begin{small}

\bibitem{II_Br} F.\,Browder:
\textit{Functional Analysis and Partial Differential Equations I}.
Math. Annalen 138, pp.\,55-79, 1959.

\bibitem{II_RS1} M.\,Reed, B.\,Simon:
\textit{Modern Methods of Mathematical Physics I: {\FA}},
Academic Press Inc., 1980.

\bibitem{II_Kr} S.\,Krein:
\textit{Linear Equations in Banach Spaces}.
Birk\-h\"{a}user, 1982.

\bibitem{II_Yo} K.\,Yosida:
\textit{Functional Analysis}.
Springer, 1980.

\end{small}
\end{thebibliography}

\newpage
\thispagestyle{empty}
\fancyhead[C]{\small{III. Adjoint Spaces and Operators}}
%------------------------------------------------------------------------------%
\part{Adjoint Spaces and Operators}                                            %
%------------------------------------------------------------------------------%
It is often useful to combine the study of a topological or normed linear space
$\banX$ with that of its \textit{dual space} $\banX^*$,
which is nothing but the set of continuous and $\F$-valued linear mappings
endowed with a linear-algebraic structure.
Dual spaces are sometimes of interest in its own right,
since they may constitute extensions of the original space $\banX$,
which has many helpful applications in higher analysis.

\section{Continuous Linear Functionals}
\section{Weak Topologies, Bidual Spaces and Reflexivity}
\section{Representation Theorems for Adjoint Spaces}
\section{Spaces of Distributions}
\section{Adjoint Operators in Locally Convex Spaces}
\section{Adjoint Operators in Banach Spaces}
\section{Hilbert Adjoints and Normal Operators}

% BIBLIOGRAPHY
%------------------------------------------------------------------------------%
\vspace{5mm}
\begin{thebibliography}{99}
\begin{small}

\bibitem{III_Ka} T.\,Kato:
\textit{{Perturbation theory of linear operators}}.
Springer, Grund\-lehren der mathematischen Wissenschaften, Band 132, 1966.

\bibitem{III_RN} F.\,Riesz, B.\,Sz.-Nagy:
\textit{{Functional Analysis}}.
Springer, 1956.

\bibitem{III_Yo} K.\,Yosida:
\textit{Functional Analysis}.
Springer, Grund\-lehren der mathematischen Wissenschaften, Band 123, 1965.

\end{small}
\end{thebibliography}

\newpage
\thispagestyle{empty}
\fancyhead[C]{\small{IV. Normally Solvable Operators}}
%------------------------------------------------------------------------------%
\part{Normally Solvable Operators}                                             %
%------------------------------------------------------------------------------%
Given a densely defined linear operator $\linA$ acting between Banach spaces
$\banX$ and\,$\banY$,
the orthogonality relations in part III can be improved
if the ranges of $\linA$ and $\linA^*$ are \textit{closed} subspaces
of $\banY$ and\,$\banX^*$,
so in these cases we actually would have
\begin{equation}\label{V_ORTHOGONALITY_RELATIONS}
{\range{\linA} } \= ^{\bot}\kernel{\linA^*}
\quad \and \quad
{\range{\linA^*} } \= \kernel{\linA}^{\bot}.
\end{equation}
One also says that
a linear operator $\linA$ is \textit{normally solvable}
if it is densely defined and
$\range{\linA} \eq \, ^{\bot}\kernel{\linA^*}$ holds,
which means that $\linA$ has closed range.

\section{Banach's Closed Range Theorem}
\section{Operators Bounded Below}
\section{The Fredholm Index}
\section{Perturbations of Semi-Fredholm Operators}
\section{Basic Theory on Fredholm Operators}
\section{Fredholm's Alternative Theorem}
\section{Concepts of Solvability for Linear Equations}

% BIBLIOGRAPHY
%------------------------------------------------------------------------------%
\vspace{5mm}
\begin{thebibliography}{99}
\begin{small}

\bibitem{IV_Br} F.\,Browder:
\textit{Functional Analysis and Partial Differential Equations I}.
Math. Annalen 138, pp.\,55-79, 1959.

\bibitem{IV_Kr} S.\,Krein:
\textit{Linear Equations in Banach Spaces}.
Birk\-h\"{a}user, 1982.

\bibitem{IV_Wl} J.\,Wloka:
\textit{Partial Differential Equations}.
Cambridge University Press, 1987.

\bibitem{IV_Yo} K.\,Yosida:
\textit{Functional Analysis}.
Springer, Grund\-lehren der mathematischen Wissenschaften, Band 123, 1965.

\end{small}
\end{thebibliography}

\newpage
\thispagestyle{empty}
\fancyhead[C]{\small{V. Calculus of Vector-Valued Functions}}
%------------------------------------------------------------------------------%
\part{Calculus of Vector-Valued Functions}                                     %
%------------------------------------------------------------------------------%
In part I we have considered topological and normed linear spaces,
whose elements are real or complex-valued functions.
There are good reasons to extend the definitions of some of these spaces
to contain so-called \textit{vector-valued} functions,
which are nothing but mappings $f\:(a,b) \xto \banX$,
where $(a,b) \subset \R$ is a finite interval
and $\banX$ is a real or complex Banach space.

\subparagraph{}
The most important spaces containing such functions are
the vector-valued {\HOELDER} spaces,
the Bochner spaces and
the Bochner-Sobolev spaces.
For the latter one the notion of weak differentiability of a vector-valued
function is required.

\section{Spaces of Classical Differentiable Functions}
\section{Riemann Integration}
\section{Linear Evolution Equations}
\section{Bochner Integration}
\section{Inhomogeneous Cauchy Problems}
\section{Weak Derivatives and Bochner Sobolev-Spaces}
\section{Hilbert-Bochner-Sobolev Spaces}
\section{Holomorphic Vector-Valued Functions}

% BIBLIOGRAPHY
%------------------------------------------------------------------------------%
\vspace{5mm}
\begin{thebibliography}{99}
\begin{small}

\bibitem{V_DU} J.\,Diestel, J.\,J.\,Uhl:
\textit{Vector Measures}.
American Mathematical Soc., 1977.

\bibitem{V_EN} K.-J.\,Engel, R.\,Nagel:
\textit{A Short Course on Operator Semigroups}.
Universitext, Springer, 2006.

\bibitem{V_Yo} K.\,Yosida:
\textit{Functional Analysis}.
Springer, Grund\-lehren der mathematischen Wissenschaften, Band 123, 1965.

\end{small}
\end{thebibliography}

\newpage
\thispagestyle{empty}
\fancyhead[C]{\small{VI. General Spectral Theory}}
%------------------------------------------------------------------------------%
\part{General Spectral Theory}                                                 %
%------------------------------------------------------------------------------%
The aim in this part
is to give an outline on the spectral theory of a closed or continuous
linear operator $\linA$ acting in a complex Banach space $\banX$.

\subparagraph{}
Let us start with the remark
that each real Banach space $\banX$ can be made
into a complex Banach space $\banX_{\C}$,
and the related procedure is denoted the \textit{complexification} of $\banX$.
It can be performed as follows:
let $\banX_{\C} \eqdef \banX + i \banX \eq \{x+iy: x,y \xeof \banX\}$
and furnished with the addition and multiplication
derived from those in $\banX$.
The norm in $\banX_{\C}$ is defined according to
\[
\norm{x+iy} \, \eqdef \,
\sup \, \{\norm{x} \cos \theta + \norm{y} \sin \theta: \theta \in [0,2\pi)\},
\]
in which this space is complete.

\paragraph{}
Let $\1_X \: \banX \xto \banX$ be the identity operator in $\banX$.
In the sequel we shall consistently use the abbreviation $\linA - \zeta$
instead of $\linA -\zeta \1_X$.
Given a linear operator
$\linA \: \domain{\linA} \subset \banX \xto \banX$,
the complex number $\zeta \xeof \C$ is called a

\begin{itemize}[leftmargin=12pt]
\item[1.]
\textit{regular value} of $\linA$,
if $\linA - \zeta \: \domain{\linA} \xto \banX$ is bijective and
the inverse operator $(\linA-\zeta)^{-1} \: \banX \xto \domain{\linA}$
is continuous.
The set $\rho(\linA)$ of regular values is denoted
the \textit{resolvent set} of $\linA$;

\item[2.]
\textit{spectral value} of $\linA$,
if $\zeta \notin \rho(\linA)$.
The set $\sigma(\linA)$ of spectral values of $\linA$ is said to be
the \textit{spectrum} of $\linA$,
thus $\sigma(\linA) \eq \C \setminus \rho(\linA)$.
\end{itemize}

\subparagraph{}
Having these concepts available,
some further questions arise {\wrt} the location and shape
of the resolvent set and the spectrum within the complex plane.
Moreover,
various classes of linear operators acting in $\banX$ can be defined
by requiring specific properties of these both fundamental sets.

\section{Resolvent Sets and Resolvent Functions}
\section{Different Kinds of Spectra}
\section{Riesz-Schauder Spectral Theory}
\section{Linear Operators with Pure Point Spectrum}
\section{The Riesz-Dunford Functional Calculus}
\section{Essential and Discrete Spectrum}

% BIBLIOGRAPHY
%------------------------------------------------------------------------------%
\vspace{5mm}
\begin{thebibliography}{99}
\begin{small}

\bibitem{VI_EN} K.J.\,Engel, R.\,Nagel:
\textit{A Short Course on Operator Semigroups},
Universitext, Springer Verlag, 2006.

\bibitem{VI_Ka} T.\,Kato:
\textit{{Perturbation theory of linear operators}}.
Springer, Grund\-lehren der mathematischen Wissenschaften, Band 132, 1966.

\bibitem{VI_RS1} M.\,Reed, B.\,Simon:
\textit{Modern Methods of Mathematical Physics I: {\FA}},
Academic Press Inc., 1980.

\bibitem{VI_Yo} K.\,Yosida:
\textit{Functional Analysis}.
Springer, Grund\-lehren der mathematischen Wissenschaften, Band 123, 1965.

\end{small}
\end{thebibliography}

\newpage
\thispagestyle{empty}
\fancyhead[C]{\small{VII. Commutative Banach Algebras}}
%------------------------------------------------------------------------------%
\part{Commutative Banach Algebras}                                             %
%------------------------------------------------------------------------------%
The aim in this part is to review commutative Banach algebras and
their applications to other topics in functional analysis and operator theory.
They are of own interest due to the celebrated \textit{Gelfand transformation}.
This presence of this fundamental mapping lead us
on the one hand
  to the \textit{Fourier transformation} in certain Lebesgue spaces and
on the other hand
  to the continuous functional calculus for \textit{normal} linear operators
  acting in Hilbert space.

\section{Normed Algebras}
\section{Resolvent and Spectrum of Banach Algebra Elements}
\section{Gelfand Representation}
\section{Involutional Algebras and C*-Algebras}
\section{The Gelfand-Naimark Theorem}
\section{Topological and Locally Compact Abelian Groups}
\section{Abstract Fourier Transformation}
\section{Fourier Transformation in Euclidean Spaces}

% BIBLIOGRAPHY
%------------------------------------------------------------------------------%
\vspace{5mm}
\begin{thebibliography}{99}
\begin{small}

\bibitem{VII_Ka} E.\,Kaniuth:
\textit{A Course in Commutative Banach Algebras},
Springer GTM 246, 2009.

\bibitem{VII_Ru} W.\,Rudin:
\textit{Fourier Analysis on Groups},
Wiley-Interscience, 1962.

\bibitem{VII_SW} E.\,Stein, G.\,Weiss:
\textit{Introduction to Fourier Analysis on Euclidean Spaces},
Princeton Mathematical Series 32, Princeton University Press, 1971.

\end{small}
\end{thebibliography}

\newpage
\thispagestyle{empty}
\fancyhead[C]{\small{VIII. Linear Operators in Hilbert Space}}
%------------------------------------------------------------------------------%
\part{Linear Operators in Hilbert Space}                                       %
%------------------------------------------------------------------------------%
The presence of the inner product in a Hilbert space can be used to
analyze linear operators with respect to this additional feature.
For example,
a linear operator $\linA$ with domain of definition $\domain{\linA}$ acting
in a Hilbert space $\hilH$ is called \textit{normal},
if $\linA \linA^* \eq \linA^* \linA$ holds,
where $\linA^*$ is the Hilbert adjoint, an operator acting in $\hilH$ as well,
which is well-defined whenever $\linA$ is densely defined.

\subparagraph{}
Important classes of normal operators are the unitary, symmetric and
self-adjoint ones.
Especially the symmetric and self-adjoint linear operators has been
extensively studied and give rise to a very complete theory.

\section{Numerical Range and External Field of Regularity}
\section{Unitary Operators and Isometries}
\section{Symmetric Linear Operators}
\section{Self-Adjoint Operators}
\section{Essentially Self-Adjointness}
\section{The von\,Neumann Extension Theory}
\section{Perturbations of Self-Adjoint Operators}

% BIBLIOGRAPHY
%------------------------------------------------------------------------------%
\vspace{5mm}
\begin{thebibliography}{99}
\begin{small}

\bibitem{VIII_EE} D.\,E.\,Edmunds, W.\,D.\,Evans:
\textit{Spectral theory and differential operators}.
Oxford University Press, 1987.

\bibitem{VIII_We} J.\,Weidmann:
\textit{Linear Operators in Hilbert Spaces}.
Springer Verlag, 1980.

\end{small}
\end{thebibliography}

\newpage
\thispagestyle{empty}
\fancyhead[C]{\small{IX. Spectral Theory of Self-Adjoint Operators}}
%------------------------------------------------------------------------------%
\part{Spectral Theory of Self-Adjoint Operators}                               %
%------------------------------------------------------------------------------%
In this part we give an outline on the spectral properties of self-adjoint
linear operators.
A linear operator $\linA$ of this type is also suitable to be decomposed
into sums of projection operators,
where the sum may be mostly countable or be given by an integral representation.
Having such a decomposition available,
it is possible to define \textit{spectral operators} of the form $f(\linA)$,
where the function $f\:\sigma(\linA) \xto \C$ is at least measurable.
The mapping $f \mapsto f(\linA)$ is called a \textit{functional calculus}
for $\linA$ and provides a useful tool to solve evolution equations
governed by a \textit{positive} and self-adjoint linear operator $\linA$.

\section{Spectral Properties of Normal Operators}
\section{Properties of the Resolvent and the Spectrum}
\section{The Discrete Spectral Theorem}
\section{Projections and Projection-Valued Measures}
\section{Functional Calculus and Spectral Operators}
\section{The Spectral Theorem for Self-Adjoint Operators}
\section{Positive Self-Adjoint Operators and Cauchy Problems}

% BIBLIOGRAPHY
%------------------------------------------------------------------------------%
\vspace{5mm}
\begin{thebibliography}{99}
\begin{small}

\bibitem{IX_RN} F.\,Riesz, B.\,Sz.-Nagy:
\textit{{Functional Analysis}}.
Springer, 1956.

\bibitem{IX_Yo} K.\,Yosida:
\textit{Functional Analysis}.
Springer, Grund\-lehren der mathematischen Wissenschaften, Band 123, 1965.

\end{small}
\end{thebibliography}

\newpage
\thispagestyle{empty}
\fancyhead[C]{\small{X. Semigroups of Bounded Linear Operators}}
%------------------------------------------------------------------------------%
\part{Semigroups of Bounded Linear Operators}                                  %
%------------------------------------------------------------------------------%
Our aim in this part
is to sketch the basic notions of \textit{semigroup theory},
which forms one of the most important topics within functional analysis
{\wrt} applications to ordinary and partial differential equations.

\subparagraph{}
The notion of \textit{semigroup} has been introduced
by \textsc{J.\,A.\,S\'{e}guier} in 1904.
Later some applications to partial differential equations were detected,
but the first systematic and comprehensive treatement of semigroup theory
was published in the pioneering monograph \cite{X_HP}
by \textsc{E.\,Hille} and \textsc{R.\,Phillips}.
Nowadays semigroup theory forms a powerful method
in solving linear evolution problems
with a multiplicity of applications to partial differential equations
governed by linear-elliptic operators.

\section{Strongly Continuous Semigroups}
\section{Infinitesimal Generators}
\section{Accretive Operators and Contractive Semigroups}
\section{Sectorial Operators and Analytic Semigroups}
\section{Skew-Adjoint Operators and Unitary Groups}
\section{Well-Posedness of Second-Order Cauchy Problems}
\section{The Cosine Family Approach}

% BIBLIOGRAPHY
%------------------------------------------------------------------------------%
\vspace{5mm}
\begin{thebibliography}{99}
\begin{small}

\bibitem{X_EN}  K.-J.\,Engel, R.\,Nagel:
\textit{A Short Course on Operator Semigroups}.
Universitext, Springer, 2006.

\bibitem{X_Go} J.\,Goldstein:
\textit{Semigroups of Linear Operators and Applications}.
Oxford University Press, 1985.

\bibitem{X_Ha} M.\,Haase:
\textit{The Functional Calculus for Sectorial Operators}.
Birk\-h\"{a}user, 2006.

\bibitem{X_HP} E.\,Hille, R.\,S.\,Phillips:
\textit{Functional Analysis and Semigroups}.
AMS Vol.\,31, 1948.

\bibitem{X_Lu} A.\,Lunardi:
\textit{Analytic Semigroups and Optimal Regularity in Parabolic Problems}.
Birk\-h\"{a}user, 1995.

\bibitem{X_Pa} A.\,Pazy:
\textit{Semigroups of Linear Operators
        and Applications to Partial Differential Equations}.
Applied Mathematical Sciences Vol.\,44, Springer, 1983.

\end{small}
\end{thebibliography}

\newpage
\thispagestyle{empty}
\fancyhead[C]{\small{XI. Coercive Sesquilinear Forms}}
%------------------------------------------------------------------------------%
\part{Coercive Sesquilinear Forms}                                             %
%------------------------------------------------------------------------------%
This part deals with the general theory of coercive sesquilinear forms
and its applications to stationary and evolutional linear equations.
Sesquilinear forms are mostly efficient
when studying abstract problems in the Hilbert space setting.

A further goal is to indicate the correspondence of
sesquilinear forms and linear operators.
In fact,
each sesquilinear form gives rise to a couple of linear operators,
and reversely,
by taking the inner products $\scalar{\linA x}{y}$,
every linear operator $\linA$ defines a sesquilinear form
on its domain of definition $\domain{\linA}$.

\subparagraph{}
The totality of applications of the presented theory
to the boundary and inital-boundary value problems
of linear partial differential equations
are often denoted as \textit{form methods} or as \textit{Hilbert space methods}.

\section{Generalities on Sesquilinear Forms}
\section{Coercive Forms and Variational Equations}
\section{Triples of Hilbert Spaces}
\section{Linear Operators Associated with a H-Elliptic Form}
\section{Sectorial Forms and m-Sectorial Operators}
\section{The Friedrichs Extension Theorem}
\section{Variational Evolution Equations}

% BIBLIOGRAPHY
%------------------------------------------------------------------------------%
\vspace{5mm}
\begin{thebibliography}{99}
\begin{small}

\bibitem{XI_AT} W.\,Arendt, A.F.M. ter Elst:
\textit{From Forms to Semigroups}.
Operator Theory: Advances and Applications, Vol.\,221, pp.\,47-69.

\bibitem{XI_Ka} T.\,Kato:
\textit{Perturbation Theory for Linear Operators}, Springer Verlag,
Grundlehren der mathematischen Wissenschaften, Band 132, 1995.

\bibitem{XI_Li} J.-L.\,Lions:
\textit{\'{E}quations Diff\'{e}rentielles Op\'{e}rationelles et Probl\`{e}mes aux Limites}.
Springer,
Grund\-lehren der mathematischen Wissenschaften, Band 111, 1961.

\end{small}
\end{thebibliography}

\newpage
\thispagestyle{empty}
\fancyhead[C]{\small{XII. Applications to Differential Operators}}
%------------------------------------------------------------------------------%
\part{Applications to Differential Operators}                                  %
%------------------------------------------------------------------------------%
This final part of this text presents applications of functional analysis
to linear differential operators and some standard problems
related with them.

\subparagraph{} %- multi-index
Remember first the notation of \textit{multi-index},
which is nothing but a $d$-tuple $\alpha \eq (\alpha_1,\ldots,\alpha_d)$
with values $\alpha_i \xeof \N_0$
and with modulus $\abs{\alpha} := \alpha_1 + \ldots + \alpha_d$.
Multi-indices are used to shorten multiple partial derivatives
$
\PA f \eqdef \partial_1^{\alpha_1}
             \partial_2^{\alpha_2} \ldots
             \partial_d^{\alpha_d} f
$
of a sufficiently often differentiable and complex-valued function $f$,
but also to denote the multiple powers
$\xi^{\alpha} := \xi_1^{\alpha_1}
\cdot \xi_2^{\alpha_2} \cdot \ldots \cdot \xi_d^{\alpha_d}$
of a vector $\,\xi \eq (\xi_1,\xi_2,\ldots,\xi_d) \xeof \C^d$.

\paragraph{} %- linear partial differential operators
By the usage of multi-indices,
a \textit{linear partial differential operator} $\A$ of even order $2m$
is defined by
\begin{equation}\label{XII_EVEN_ORDER_LPDO}
\A(x,\,\cdot\,) \eqdef
\sum_{\abs{\alpha} \leq 2m} c_{\alpha}(x) \, \PA
\end{equation}
for $x \xeof \clG$,
the closure of a \textit{domain} $\G$,
that is,
an open and connected subset of an Euclidean space $\Rd$.
This domain may be bounded or unbounded and even the whole of $\Rd$.

\subparagraph{} %- essentially real differential operators
We shall call $\A$ \textit{essentially real},
if the coefficients $c_{\alpha}$ of highest order $\abs{\alpha} \eq 2m$
are real-valued.
Moreover,
the expression $\A$ is said to be
an \textit{operator with constant coefficients},
if all functions $c_{\alpha}$ are constant in $\clG$.

\section{Free Space and Boundary Value Problems}
\section{Operators of Elliptic, Parabolic and Hyperbolic Type}
\section{Realizations in Banach Spaces of Functions}
\section{Closeness and Closability of Differential Operators}
\section{The Fredholm Property of Properly Elliptic Operators}
\section{Self-Adjoint Extensions}
\section{Eigenvalue Problems in Bounded Domains}
\section{Variational Boundary Value Problems}
\section{Linear Differential Operators in Distributional Spaces}
\section{Fourier Transform Methods}

% BIBLIOGRAPHY
%------------------------------------------------------------------------------%
\vspace{5mm}
\begin{thebibliography}{99}
\begin{small}

\bibitem{XII_Br2} F.\,Browder:
\textit{On the Spectral Theory of Elliptic Differential Operators I}.
Math. Annalen 142, pp.\,22-130, 1961.

\bibitem{XII_GT}  D.\,Gilbarg, N.S.\,Trudinger:
\textit{Elliptic Partial Differential Operators of Second Order}.
Springer, 2001.

\bibitem{XII_Pa} A.\,Pazy:
\textit{Semigroups of Linear Operators
        and Applications to Partial Differential Equations}.
Applied Mathematical Sciences Vol.\,44, Springer, 1983.

\bibitem{XII_RR} M.\,Renardy, R.\,Rogers:
\textit{An Introduction to Partial Differential Equations}.
Texts in Applied Mathematics 13, Springer, 2004.

\bibitem{XII_Yo} K.\,Yosida:
\textit{Functional Analysis}.
Springer, Grund\-lehren der mathematischen Wissenschaften, Band 123, 1965.

\end{small}
\end{thebibliography}

%------------------------------------------------------------------------------%
%                               END OF DOCUMENT                                %
%------------------------------------------------------------------------------%
\end{document}

   I. TOPOLOGICAL AND NORMED LINEAR SPACES
      ------------------------------------
   1. Generalities on Topological Linear Spaces
   2. Seminorms and Locally Convex Spaces
   3. Normed Linear Spaces
   4. Strictly and Uniformly Convex Spaces
   5. Banach and Hilbert Spaces
   6. Spaces of Scalar-Valued Functions and Measures
   7. Hierarchy of Topological and Normed Linear Spaces

  II. LINEAR OPERATORS IN BANACH SPACES
      ---------------------------------
   1. Fundamentals of Operator Theory
   2. Continuous Linear Operators
   3. Compact Linear Operators
   4. Closed and Closable Linear Operators
   5. The Closed Graph Theorem

 III. ADJOINT SPACES AND OPERATORS
      ----------------------------
   1. Continuous Linear Functionals
   2. Weak Topologies, Bidual Spaces and Reflexivity
   3. Representation Theorems for Adjoint Spaces
   4. Spaces of Distributions
   5. Adjoint Operators in Locally Convex Spaces
   6. Adjoint Operators in Banach Spaces
   7. Hilbert Adjoints and Normal Operators

  IV. NORMALLY SOLVABLE OPERATORS
      ---------------------------
   1. Banach's Closed Range Theorem
   2. Operators Bounded Below
   3. The Fredholm Index
   4. Perturbations of Semi-Fredholm Operators
   5. Basic Theory on Fredholm Operators
   6. Fredholm's Alternative Theorem
   7. Concepts of Solvability for Linear Equations

   V. CALCULUS OF VECTOR-VALUED FUNCTIONS
      -----------------------------------
   1. Spaces of Classical Differentiable Functions
   2. Riemann Integration
   3. Linear Evolution Equations
   4. Bochner Integration
   5. Inhomogeneous Cauchy Problems
   6. Weak Derivatives and Bochner Sobolev-Spaces
   7. Hilbert-Bochner-Sobolev Spaces
   8. Holomorphic Vector-Valued Functions

  VI. GENERAL SPECTRAL THEORY
      -----------------------
   1. Resolvent Sets and Resolvent Functions
   2. Different Kinds of Spectra
   3. Riesz-Schauder Spectral Theory
   4. Linear Operators with Pure Point Spectrum
   5. The Riesz-Dunford Functional Calculus
   6. Essential and Discrete Spectrum

 VII. COMMUTATIVE BANACH ALGEBRAS
      ---------------------------
   1. Normed Algebras
   2. Resolvent and Spectrum of Banach Algebra Elements
   3. Gelfand Representation
   4. Involutional Algebras and C*-Algebras
   5. The Gelfand-Naimark Theorem
   6. Topological and Locally Compact Abelian Groups
   7. Abstract Fourier Transformation
   8. Fourier Transformation in Euclidean Spaces

VIII. LINEAR OPERATORS IN HILBERT SPACE
      ---------------------------------
   1. Numerical Range and External Field of Regularity
   2. Unitary Operators and Isometries
   3. Symmetric Linear Operators
   4. Self-Adjoint Operators
   5. Essentially Self-Adjointness
   6. The von Neumann Extension Theory
   7. Perturbations of Self-Adjoint Operators

  IX. SPECTRAL THEORY OF SELF-ADJOINT OPERATORS
      -----------------------------------------
   1. Spectral Properties of Normal Operators
   2. Properties of the Resolvent and the Spectrum
   3. The Discrete Spectral Theorem
   4. Projections and Projection-Valued Measures
   5. Functional Calculus and Spectral Operators
   6. The Spectral Theorem for Self-Adjoint Operators
   7. Positive Self-Adjoint Operators and Cauchy Problems

   X. SEMIGROUPS OF BOUNDED LINEAR OPERATORS
      --------------------------------------
   1. Strongly Continuous Semigroups
   2. Infinitesimal Generators
   3. Accretive Operators and Contractive Semigroups
   4. Sectorial Operators and Analytic Semigroups
   5. Skew-Adjoint Operators and Unitary Groups
   6. Well-Posedness of Second-Order Cauchy Problems
   7. The Cosine Family Approach

  XI. COERCIVE SESQUILINEAR FORMS
      ---------------------------
   1. Generalities on Sesquilinear Forms
   2. Coercive Forms and Variational Equations
   3. Triples of Hilbert Spaces
   4. Linear Operators Associated with a H-Elliptic Form
   5. Sectorial Forms and m-Sectorial Operators
   6. The Friedrichs Extension Theorem
   7. Variational Evolution Equations

 XII. APPLICATIONS TO DIFFERENTIAL OPERATORS
      --------------------------------------
   1. Free Space and Boundary Value Problems
   2. Operators of Elliptic, Parabolic and Hyperbolic Type
   3. Realizations in Banach Spaces of Functions
   4. Closeness and Closability of Differential Operators
   5. The Fredholm Property of Properly Elliptic Operators
   6. Self-Adjoint Extensions
   7. Eigenvalue Problems in Bounded Domains
   8. Variational Boundary Value Problems
   9. Linear Differential Operators in Distributional Spaces
  10. Fourier Transform Methods
